\documentclass[10pt,a4paper]{article}

\usepackage{lmodern}
\usepackage{microtype}
\usepackage[a4paper,left=2.75cm,right=2.75cm,top=2.75cm,bottom=2.75cm]{geometry}
\usepackage{amsmath,amssymb,amsthm,mathtools}
\usepackage{booktabs}
\usepackage[colorlinks=true,linkcolor=black,citecolor=black,urlcolor=blue]{hyperref}

\numberwithin{equation}{section}

\theoremstyle{plain}
\newtheorem{theorem}{Theorem}[section]
\newtheorem{lemma}[theorem]{Lemma}
\newtheorem{proposition}[theorem]{Proposition}
\newtheorem{corollary}[theorem]{Corollary}

\theoremstyle{definition}
\newtheorem{definition}[theorem]{Definition}

\theoremstyle{remark}
\newtheorem*{remark}{Remark}

\DeclareMathOperator{\Tr}{Tr}
\DeclareMathOperator{\supp}{supp}
\DeclareMathOperator{\diag}{diag}
\DeclareMathOperator{\QJS}{QJS}
\DeclareMathOperator{\BSC}{BSC}
\DeclareMathOperator{\BEC}{BEC}
\DeclareMathOperator{\Pure}{Pure}
\DeclareMathOperator{\Unif}{Unif}
\DeclareMathOperator{\id}{id}

\newcommand{\cH}{\mathcal{H}}
\newcommand{\cD}{\mathcal{D}}
\newcommand{\cB}{\mathcal{B}}
\newcommand{\cZ}{\mathcal{Z}}
\newcommand{\cL}{\mathcal{L}}
\newcommand{\cN}{\mathcal{N}}
\newcommand{\cE}{\mathcal{E}}
\newcommand{\cF}{\mathcal{F}}
\newcommand{\cG}{\mathcal{G}}
\newcommand{\cA}{\mathcal{A}}
\newcommand{\cP}{\mathcal{P}}
\newcommand{\cY}{\mathcal{Y}}
\newcommand{\bC}{\mathbb{C}}
\newcommand{\bE}{\mathbb{E}}
\newcommand{\bP}{\mathbb{P}}
\newcommand{\ket}[1]{\lvert #1 \rangle}
\newcommand{\bra}[1]{\langle #1 \rvert}
\newcommand{\ketbra}[1]{\lvert #1 \rangle\langle #1 \rvert}
\newcommand{\braket}[2]{\langle #1 \vert #2 \rangle}

\title{Polynomial-Blocklength Quantum Polar Codes\\ via Nonlinear Dephasing Contraction}
\author{Prince Anand}
\date{November 2026}

\begin{document}

\maketitle

\begin{quote}\small
\noindent\textbf{Abstract.} We prove that polar codes achieve the symmetric capacity of every finite-dimensional binary-input classical--quantum channel at blocklength $N = O((1/\varepsilon)^{\mu})$ with average block-error probability at most $2^{-N^{0.49}}$ for a universal constant $\mu < \infty$. Starting from Frenkel's integral formula for quantum relative entropy, we derive a quantum Jeffreys--Jensen--Shannon inequality and integrate it along the qubit dephasing semigroup to obtain a nonlinear strong data processing inequality for $\cZ_p \otimes \id$. This proves the Hirche--Reeb Quantum Mrs.\ Gerber's Lemma whenever one combined channel is pure-state or classical, and yields $H(W^-) - H(W) \ge \theta_0 H(W)(1 - H(W))$ ($\theta_0 > 0.05673$) for all binary-input classical--quantum channels.
\end{quote}

\section{Introduction}

\subsection{Background and Motivation}

A central objective in both classical and quantum information theory is the design of structured coding schemes that approach channel capacity with low complexity and well-controlled finite-blocklength performance. For classical channels, Ar{\i}kan~\cite{Arikan2009} introduced polar codes and proved that they achieve the symmetric capacity of every binary-input discrete memoryless channel (B-DMC) with $O(N \log N)$ encoding and successive-cancellation decoding complexity. Given $N = 2^n$ independent copies of a binary-input channel $W$, Ar{\i}kan's recursive CNOT transform synthesizes $N$ virtual bit-channels $\{W_n^{(i)}\}_{i=1}^{N}$ whose conditional entropies polarize toward the extremes $0$ and $1$: a fraction approaching the symmetric capacity $I(W)$ becomes nearly noiseless and carries information bits, while the complementary fraction becomes nearly completely noisy and is frozen to values known to both sender and receiver. Soon thereafter, Ar{\i}kan and Telatar~\cite{ArikanTelatar2009} determined the asymptotic error-exponent behavior of this construction, showing that for any fixed rate $R < I(W)$ and any $\beta < 1/2$, the block-error probability under successive-cancellation decoding decays as $o(2^{-N^{\beta}})$.

In the quantum setting, the classical capacity of a classical--quantum (c-q) channel is characterized by the Holevo--Schumacher--Westmoreland coding theorem~\cite{Holevo1998,SchumacherWestmoreland1997}. Wilde and Guha~\cite{WildeGuha2013} extended Ar{\i}kan's construction to arbitrary binary-input c-q channels $W : \{0,1\} \to \cD(\cH)$ mapping classical input bits to density operators on a finite-dimensional Hilbert space $\cH$. Using a quantum successive-cancellation decoder built from Helstrom measurements~\cite{Helstrom1976} and analyzed via Sen's noncommutative union bound~\cite{Sen2012} (see also the sharper sequential union bounds of Gao~\cite{Gao2015} and Oskouei, Mancini, and Wilde~\cite{OskoueiManciniWilde2019}), they established asymptotic polarization and proved that for any fixed rate $R < I(W)$ and any $\beta < 1/2$, the block-error probability decays as $o(2^{-\frac{1}{2} N^{\beta}})$. Subsequent developments extended polar coding to degradable quantum channels~\cite{WildeGuha2013Degradable}, private classical and entanglement-assisted quantum communication over arbitrary qubit-input channels~\cite{RenesDupuisRenner2012,RenesWilde2014}, unassisted quantum communication via chained polar blocks~\cite{RenesSutterDupuisRenner2015}, and general non-binary c-q channels and multiple-access channels~\cite{NasserRenes2018}. Because these polarization proofs—like Ar{\i}kan's original argument—rely on the martingale convergence theorem, they establish asymptotic achievability at any fixed rate below capacity, yet provide no quantitative bound on how large the blocklength $N$ must be in order to operate within a prescribed additive gap $\varepsilon > 0$ of the symmetric Holevo capacity $I(W)$.

In classical coding theory, this finite-blocklength question—often phrased as determining the scaling exponent $\mu$ governing the tradeoff $N = O((1/\varepsilon)^{\mu})$ between blocklength and gap to capacity—was resolved by Guruswami and Xia~\cite{GuruswamiXia2015} and independently by Hassani, Alishahi, and Urbanke~\cite{HassaniAlishahiUrbanke2014}, with subsequent refinements on the scaling exponent and unified moderate-deviations regimes by Goldin and Burshtein~\cite{GoldinBurshtein2014} and Mondelli, Hassani, and Urbanke~\cite{MondelliHassaniUrbanke2016}. These works established that there exists a universal constant $\mu < \infty$ such that classical polar codes achieve rates $R \ge I(W) - \varepsilon$ with exponentially small block-error probability at polynomial blocklength $N = O((1/\varepsilon)^{\mu})$. The analytic engine behind this polynomial scaling is a quantitative one-step entropy-splitting lower bound that behaves linearly at both endpoints $H(W) \to 0$ and $H(W) \to 1$:
\begin{equation}\label{eq:linear-splitting-intro}
H(W^-) - H(W) = H(W) - H(W^+) \;\ge\; \theta\, H(W)\bigl(1 - H(W)\bigr),
\end{equation}
valid for every binary-input classical channel with a universal constant $\theta > 0$. Once~\eqref{eq:linear-splitting-intro} is available, the polarization potential $T(W) = \sqrt{H(W)(1 - H(W))}$ contracts by a uniform multiplicative factor at every step of the polar tree, allowing the fraction of unpolarized channels to be bounded after $O(\log(1/\varepsilon))$ levels without appealing to asymptotic martingale convergence.

For classical channels, \eqref{eq:linear-splitting-intro} is a direct consequence of Wyner and Ziv's Mrs.\ Gerber's Lemma~\cite{WynerZiv1973}, which states that for two independent binary random variables $X_1, X_2$ with classical side information $Y_1, Y_2$,
\begin{equation}\label{eq:classical-mgl}
H(X_1 \oplus X_2 \mid Y_1 Y_2) \;\ge\; h_2\!\left(h_2^{-1}(H(X_1 \mid Y_1)) * h_2^{-1}(H(X_2 \mid Y_2))\right),
\end{equation}
where $h_2$ is the binary entropy function and $*$ denotes binary convolution. The proof of~\eqref{eq:classical-mgl} relies on two structural features of classical conditioning: the convexity of $(u,v) \mapsto h_2(h_2^{-1}(u) * h_2^{-1}(v))$ in each argument, and the decomposition $H(X \mid Y) = \sum_y \bP(Y=y)\, H(X \mid Y=y)$, which expresses conditional Shannon entropy as an average of unconditioned binary entropies and therefore reduces the general inequality to binary symmetric channels via Jensen's inequality.

When the side information is quantum, both pillars of this argument fail. Conditional von Neumann entropy $H(X \mid B)$ cannot in general be decomposed as a convex combination of unconditioned entropies, and in fact pure-state c-q channels can strictly violate~\eqref{eq:classical-mgl} (see~\cite{HircheReeb2018}). Hirche and Reeb~\cite{HircheReeb2018} initiated a systematic study of binary information combining with quantum side information by combining Fawzi and Renner's fidelity-of-recovery inequality~\cite{FawziRenner2015}, a fidelity--entropy relation, and Renes' channel duality~\cite{Renes2018}. This yielded a nontrivial lower bound on $H(W^-) - H(W)$ for all binary-input c-q channels. Near the boundary $\delta := \min(H(W), 1 - H(W)) \to 0$, however, their bound scales as $\Omega(\delta/\log(1/\delta))$ rather than $\Omega(\delta)$. Because the contraction factor of the polarization potential degrades logarithmically near the endpoints, their finite-blocklength analysis yielded only a sub-exponential blocklength bound $N = O((1/\varepsilon)^{\mu \log(1/\varepsilon)})$~\cite[Theorem~VIII.4]{HircheReeb2018}.

Moreover, Hirche and Reeb demonstrated~\cite[Remark~VI.4]{HircheReeb2018} that this logarithmic loss is structural for fidelity-of-recovery methods: even an idealized Petz-map recovery bound loses a factor of $\log(1/\delta)$ as $H(W) \to 1$, because fidelity varies quadratically under small state perturbations whereas von Neumann entropy varies logarithmically. They formulated a duality-symmetrized counterpart of~\eqref{eq:classical-mgl} as the conjectured optimal Quantum Mrs.\ Gerber's Lemma~\cite[Conjecture~VII.1]{HircheReeb2018}, and observed~\cite[Remark~VIII.5]{HircheReeb2018} that any linear-scaling lower bound of the form~\eqref{eq:linear-splitting-intro} would suffice to establish polynomial-blocklength classical--quantum polar codes. (The complementary erasure-channel upper bound, \cite[Conjecture~VII.2]{HircheReeb2018}, was recently addressed in~\cite{He2026}.)

\subsection{Main Results}

In this paper we prove that classical--quantum polar codes achieve symmetric capacity at polynomial blocklength for all finite-dimensional binary-input c-q channels, with a universal scaling exponent $\mu$ independent of both the channel and the output Hilbert-space dimension. We also prove Conjecture~VII.1 of Hirche and Reeb~\cite{HircheReeb2018} whenever one of the two combined channels is pure-state or classical. Rather than appealing to recovery-map inequalities, our approach derives a comparison inequality between symmetrized quantum relative entropy and quantum Jensen--Shannon divergence from Frenkel's integral representation~\cite{Frenkel2023} and integrates it along the continuous-time qubit dephasing semigroup.

Throughout the paper, all logarithms are base $2$, entropies and divergences are measured in bits, $h_2(p) := -p \log_2 p - (1-p)\log_2(1-p)$ denotes the binary entropy with inverse $h_2^{-1} : [0,1] \to [0,1/2]$, and $a * b := a(1-b) + (1-a)b$ denotes binary convolution. We define
\[
g_c(u,v) := h_2\!\bigl(h_2^{-1}(u) * h_2^{-1}(v)\bigr), \qquad
g_q(u,v) := \min\bigl\{g_c(u,v),\; u + v - 1 + g_c(1-u, 1-v)\bigr\},
\]
for $(u,v) \in [0,1]^2$, and $\Delta(u,p) := h_2(h_2^{-1}(u) * p) - h_2(p)$ for $u \in [0,1]$ and $p \in [0,1/2]$. Formal definitions of all entropic quantities, divergences, and polar channel combinations are collected in Section~\ref{sec:preliminaries}.

Our principal coding result is the following polynomial-blocklength theorem.

\begin{theorem}[Polynomial-Blocklength Classical--Quantum Polar Codes]\label{thm:main-coding}
There exist absolute constants $\mu, C_0 \in (0,\infty)$ such that for every binary-input c-q channel $W : \{0,1\} \to \cD(\cH)$ on any finite-dimensional Hilbert space $\cH$, every $\varepsilon \in (0,1)$, and every power-of-two blocklength
\begin{equation}\label{eq:blocklength-bound}
N = 2^n \;\ge\; C_0 \left(\frac{1}{\varepsilon}\right)^{\mu},
\end{equation}
there exists a $G_N$-coset classical--quantum polar code of blocklength $N$ and rate $R \ge I(W) - \varepsilon$ whose average quantum successive-cancellation block-error probability (under uniformly distributed messages) satisfies $P_e(N,R) \le 2^{-N^{0.49}}$.
\end{theorem}

Theorem~\ref{thm:main-coding} is driven by our second main result, which establishes Conjecture~VII.1 of~\cite{HircheReeb2018} for all hybrid channel pairs and yields a linear-scaling Quantum Mrs.\ Gerber bound for arbitrary pairs of binary-input c-q channels.

\begin{theorem}[Quantum Mrs.\ Gerber Bounds]\label{thm:qmgl}
Let $W_1 : x \mapsto \rho_x \in \cD(\cH_1)$ and $W_2 : x \mapsto \sigma_x \in \cD(\cH_2)$ be independent binary-input c-q channels with conditional entropies $H_1 = H(W_1)$ and $H_2 = H(W_2)$.
\begin{enumerate}
\item[(a)] \emph{Hybrid pairs.} If $W_2$ is a pure-state channel, then
\begin{equation}\label{eq:qmgl-pure}
H(W_1 \boxtimes W_2) \;\ge\; H_1 + H_2 - 1 + g_c(1 - H_1, 1 - H_2) \;\ge\; g_q(H_1, H_2).
\end{equation}
If $W_2$ is a classical binary-input channel, then
\begin{equation}\label{eq:qmgl-classical}
H(W_1 \boxtimes W_2) \;\ge\; g_c(H_1, H_2) \;\ge\; g_q(H_1, H_2).
\end{equation}
\item[(b)] \emph{General channels.} For arbitrary $W_1, W_2$,
\begin{equation}\label{eq:four-branch}
\begin{aligned}
H(W_1 \boxtimes W_2) \;\ge\; \max\Bigl\{
&H_1 + \Delta(1 - H_1, 1 - 2^{H_2 - 1}),\;
H_2 + \Delta(1 - H_2, 1 - 2^{H_1 - 1}),\\
&H_2 + \Delta(H_1, 1 - 2^{-H_2}),\;
H_1 + \Delta(H_2, 1 - 2^{-H_1})
\Bigr\}.
\end{aligned}
\end{equation}
In particular, for every binary-input c-q channel $W$ with $H = H(W) \in [0,1]$,
\begin{equation}\label{eq:linear-splitting}
H(W^-) - H(W) = H(W) - H(W^+) \;\ge\; \theta_0\, H(W)\bigl(1 - H(W)\bigr),
\end{equation}
where $\theta_0 := \frac{32}{3}\, h_2^{-1}(1/4)\bigl(1 - h_2(1 - 2^{-1/2})\bigr) = 0.05673348\ldots > 0.05673$, with endpoint scaling
\begin{equation}\label{eq:endpoint-scaling}
\lim_{\delta \to 0^+} \frac{\Delta(\delta, 1 - 2^{-\delta})}{\delta} \;=\; 1.
\end{equation}
Consequently, the polarization potential $T(W) := \sqrt{H(W)(1 - H(W))}$ contracts uniformly at every step:
\begin{equation}\label{eq:potential-contraction}
\frac{1}{2}\bigl(T(W^+) + T(W^-)\bigr) \;\le\; \Lambda_0\, T(W), \qquad
\Lambda_0 := \sqrt{1 - \frac{\theta_0^2}{4}} = 0.99959758\ldots < 0.99960.
\end{equation}
\end{enumerate}
\end{theorem}

The limit~\eqref{eq:endpoint-scaling} illustrates the mechanism by which the logarithmic barrier of~\cite{HircheReeb2018} is overcome. As $\delta = 1 - H \to 0^+$, the step size $p_1 = h_2^{-1}(\delta) \sim \delta/\log_2(1/\delta)$ carries an inverse-logarithmic factor, while the derivative of $h_2$ evaluated at $p_2 = 1 - 2^{-\delta} \sim (\ln 2)\delta$ contributes $h_2'(p_2) \sim \log_2(1/\delta)$. The two logarithmic factors cancel, restoring linear scaling in $\delta$ with asymptotic coefficient~$1$.

Both parts of Theorem~\ref{thm:qmgl} follow from a stabilized, dimension-independent nonlinear strong data processing inequality (SDPI) for the qubit dephasing channel. Nonlinear SDPIs—which bound output divergence as a concave function of input divergence rather than by a single linear contraction coefficient—play a central role in classical and quantum information theory~\cite{PolyanskiyWu2017,Raginsky2016,HircheRouzeFranca2022}. In the resource theory of quantum coherence~\cite{BaumgratzCramerPlenio2014}, the quantity $C(p)$ below is the relative entropy of coherence on qubit $A$ in the presence of an arbitrary quantum memory $B$.

\begin{theorem}[Nonlinear SDPI for Qubit Dephasing]\label{thm:dephasing-sdpi}
Let $\cZ_p(\cdot) := (1-p)(\cdot) + p\, Z(\cdot)Z$ denote the qubit dephasing channel with flip probability $p \in [0, 1/2]$. For every finite-dimensional Hilbert space $\cH_B$ and every bipartite state $\Psi_{AB} \in \cD(\bC^2 \otimes \cH_B)$, let
\[
C(p) := D\!\bigl((\cZ_p \otimes \id)(\Psi_{AB}) \,\Vert\, (\cZ_{1/2} \otimes \id)(\Psi_{AB})\bigr) \;\in\; [0,1]
\]
denote the relative entropy of coherence on qubit $A$. Then
\begin{equation}\label{eq:dephasing-sdpi}
C(p) \;\le\; 1 - h_2\!\bigl(h_2^{-1}(1 - C(0)) * p\bigr).
\end{equation}
Equality is attained for every $p \in [0, 1/2]$ by every equatorial single-qubit state $\Psi_A = \frac{1}{2}(I + xX)$ with $x \in [0,1]$.
\end{theorem}

Theorem~\ref{thm:dephasing-sdpi} shows that arbitrary entanglement between the dephased qubit $A$ and an unrestricted quantum memory $B$ cannot cause the relative entropy of coherence to contract more slowly than it does for an unentangled equatorial qubit state. In particular, the bound~\eqref{eq:dephasing-sdpi} is tight at every noise level $p \in [0,1/2]$ and independent of $\dim \cH_B$, whereas the linear contraction coefficient of $\cZ_p$ is trivial ($\eta_{\mathrm{Re}}(\cZ_p) = 1$, cf.~\cite{HircheRouzeFranca2022}) because classical diagonal states are fixed by $\cZ_p$.

\subsection{Proof Strategy and Organization of the Paper}

The proof links four areas: integral representations of quantum divergences, entropy production along quantum Markov semigroups, binary quantum information combining, and multi-stage martingale polarization.

Section~\ref{sec:preliminaries} reviews the entropic quantities, distinguishability measures, binary-input c-q channels, polar channel combinations ($W_1 \boxtimes W_2$ and $W_1 \circledast W_2$), entropy conservation, and Renes' channel duality~\cite{Renes2018}. Particular care is taken with the variable-node channel $W_1 \circledast W_2$, whose classical register $U_1$ is required for entropy conservation and channel duality on asymmetric channels.

Section~\ref{sec:qjjs-sdpi} proves Theorem~\ref{thm:dephasing-sdpi}. Starting from Frenkel's integral representation of quantum relative entropy~\cite{Frenkel2023} in its hockey-stick form~\cite{HircheTomamichel2024} (see also~\cite{BertaLamiTomamichel2024}), together with the midpoint identity $E_s(\rho \Vert \frac{\rho+\sigma}{2}) = (1 - \frac{s}{2}) E_{\frac{s}{2-s}}(\rho \Vert \sigma)$ for $s \in (0,2)$, we show that $D(\rho\Vert\sigma)$, $D(\sigma\Vert\rho)$, and $\QJS(\rho,\sigma)$ admit simultaneous representations as classical $f$-divergences against a single compactly supported probability measure $d\mu(t) = d(f'_+(t))$ determined by $f(s) = E_s(\rho\Vert\sigma)$ (Lemma~\ref{lem:hockey-stick-measure}). Integrating a perspective-symmetric scalar convexity inequality (Lemma~\ref{lem:scalar-convexity}) against $\mu$ yields the quantum Jeffreys--Jensen--Shannon inequality (Proposition~\ref{prop:qjjs}). In the spirit of entropy-production methods for quantum Markov semigroups~\cite{KastoryanoTemme2013,MullerHermesFrancaWolf2016,BeigiDattaRouze2020}, differentiating $C(p(t)) = S(\Psi_\infty) - S(\Psi_t)$ along the dephasing semigroup $\cZ_{p(t)} = e^{t\cL}$ converts Proposition~\ref{prop:qjjs} into the nonlinear differential inequality $-\frac{dC}{dt} \ge \phi(C)$. Under the reparametrization $x(t) = 1 - 2h_2^{-1}(1 - C(p(t)))$, this linearizes to $\frac{dx}{dt} \le -x(t)$ and integrates to Theorem~\ref{thm:dephasing-sdpi}.

Section~\ref{sec:qmgl} proves Theorem~\ref{thm:qmgl}. A Gram-matrix factorization (Lemma~\ref{lem:gram-dephasing}) converts the combination of an arbitrary c-q channel $W_1$ with a pure-state channel $W_2$ into local qubit dephasing on a bipartite state of initial coherence $H(W_1)$. Unitary reflection symmetry of pure-state channels then yields~\eqref{eq:qmgl-pure}, while channel duality for $\BSC(q)$ and classical posterior conditioning yield~\eqref{eq:qmgl-classical} for arbitrary (possibly asymmetric) classical $W_2$. Part~(b) is established in three steps (Propositions~\ref{prop:four-branch}--\ref{prop:potential-contraction}) via Uhlmann purification, a two-region convexity analysis of $\Delta(1 - H, 1 - 2^{H-1})$, and a square-root contraction estimate for $T(W)$.

Section~\ref{sec:polar-codes} proves Theorem~\ref{thm:main-coding}. After establishing the quantum Bhattacharyya evolution inequalities (Lemma~\ref{lem:bhattacharyya}), we track the synthetic channel process $(W_t)_{t \ge 0}$ through three consecutive stages of depths $m_1, m_2, m_3 = O(\log(1/\varepsilon))$: rough entropy polarization via~\eqref{eq:potential-contraction} (Lemma~\ref{lem:stage1}), fidelity bootstrapping below $2^{-Y_0}$ via a Doob-stopped $\sqrt{Z}$-supermartingale (Lemma~\ref{lem:stage2}), and stretched-exponential decay via a golden-ratio exponential supermartingale combined with Hoeffding's inequality~\cite{Hoeffding1963} (Lemma~\ref{lem:stage3}). Applying Wilde and Guha's successive-cancellation error bound~\cite[Proposition~4]{WildeGuha2013} completes the proof of Theorem~\ref{thm:main-coding}, and Corollary~\ref{cor:private-quantum} records the resulting polynomial-blocklength scaling for entanglement-assisted quantum and secret-key-assisted private classical polar coding via Renes and Wilde~\cite{RenesWilde2014}. Section~\ref{sec:concluding} closes with remarks on open problems.

\section{Preliminaries and Channel Duality}\label{sec:preliminaries}

This section fixes notation and collects the standard properties of quantum entropies, divergences, binary-input c-q channels, polar combinations, and channel duality used throughout the paper. For general background on quantum information theory and finite-resource entropic quantities, we refer to~\cite{Wilde2013,Tomamichel2016}.

\subsection{Entropies, Divergences, and Distinguishability Measures}

Throughout, $\cH$ denotes a finite-dimensional complex Hilbert space, $\cB(\cH)$ denotes the algebra of linear operators on $\cH$, and $\cD(\cH) := \{\rho \in \cB(\cH) : \rho \succeq 0,\, \Tr \rho = 1\}$ denotes the set of density operators (quantum states) on $\cH$. For a Hermitian operator $A$ on $\cH$, let $A_+$ and $A_-$ denote its positive and negative parts in the spectral decomposition, so that $A = A_+ - A_-$ with $A_+, A_- \succeq 0$ and $A_+ A_- = 0$. The support of a positive semidefinite operator $\rho$ is $\supp(\rho) := (\ker \rho)^{\perp}$.

For $\rho, \sigma \in \cD(\cH)$, the \emph{von Neumann entropy} of $\rho$ is
\[
S(\rho) := -\Tr(\rho \log_2 \rho) \;\in\; [0, \log_2 \dim \cH],
\]
with the convention $0 \log_2 0 := 0$. The \emph{Umegaki quantum relative entropy}~\cite{Umegaki1962} is defined by
\[
D(\rho \Vert \sigma) :=
\begin{cases}
\Tr\!\bigl[\rho(\log_2 \rho - \log_2 \sigma)\bigr], & \text{if } \supp(\rho) \subseteq \supp(\sigma),\\
+\infty, & \text{otherwise.}
\end{cases}
\]
It satisfies Klein's inequality $D(\rho \Vert \sigma) \ge 0$, with equality if and only if $\rho = \sigma$, and the data processing inequality $D(\cN(\rho) \Vert \cN(\sigma)) \le D(\rho \Vert \sigma)$ for every completely positive trace-preserving (CPTP) map $\cN$~\cite{Lindblad1975} (equivalent to strong subadditivity of von Neumann entropy~\cite{LiebRuskai1973}, and in fact valid for all positive trace-preserving maps~\cite{MullerHermesReeb2017}). The \emph{quantum Jensen--Shannon divergence}—equivalently, the Holevo information~\cite{Holevo1973} of the uniform binary ensemble $\{\tfrac{1}{2}, \rho; \tfrac{1}{2}, \sigma\}$—is
\begin{equation}\label{eq:qjs-def}
\QJS(\rho, \sigma) := S\!\left(\frac{\rho + \sigma}{2}\right) - \frac{1}{2}\bigl(S(\rho) + S(\sigma)\bigr)
= \frac{1}{2} D\!\left(\rho \,\middle\Vert\, \frac{\rho + \sigma}{2}\right) + \frac{1}{2} D\!\left(\sigma \,\middle\Vert\, \frac{\rho + \sigma}{2}\right).
\end{equation}
By concavity of von Neumann entropy and the Holevo bound~\cite{Holevo1973}, we have $0 \le \QJS(\rho, \sigma) \le 1$, with $\QJS(\rho, \sigma) = 0$ if and only if $\rho = \sigma$ and $\QJS(\rho, \sigma) = 1$ if and only if $\rho \perp \sigma$ (meaning $\supp(\rho) \perp \supp(\sigma)$). Moreover, $\QJS$ is symmetric in its arguments and invariant under simultaneous unitary conjugation: $\QJS(U\rho U^{\dagger}, U\sigma U^{\dagger}) = \QJS(\rho, \sigma)$ for every unitary $U$ on $\cH$.

Our integral representation in Section~\ref{sec:qjjs-sdpi} is built on the \emph{quantum hockey-stick divergence} of order $\gamma \ge 0$,
\begin{equation}\label{eq:hockey-stick-def}
E_{\gamma}(\rho \Vert \sigma) := \Tr\!\bigl[(\rho - \gamma \sigma)_+\bigr] = \max_{0 \preceq \Pi \preceq I} \Tr\!\bigl[\Pi(\rho - \gamma \sigma)\bigr],
\end{equation}
where the maximum ranges over all effect operators $0 \preceq \Pi \preceq I$ and is attained by the orthogonal projection onto the positive eigenspace of $\rho - \gamma \sigma$. Operationally, $E_\gamma(\rho \Vert \sigma)$ governs asymmetric binary quantum hypothesis testing with threshold $\gamma$~\cite{Helstrom1976,HircheTomamichel2024}. At $\gamma = 1$, it reduces to the \emph{trace distance}
\[
T_{\mathrm{tr}}(\rho, \sigma) := \frac{1}{2}\lVert \rho - \sigma \rVert_1 = E_1(\rho \Vert \sigma),
\]
which determines the optimal probability of distinguishing $\rho$ from $\sigma$ under uniform prior~\cite{Helstrom1976}.

Finally, the \emph{Uhlmann root fidelity}~\cite{Uhlmann1976} (also referred to as the quantum Bhattacharyya parameter in polar coding~\cite{WildeGuha2013}) is
\[
F(\rho, \sigma) := \lVert \sqrt{\rho}\sqrt{\sigma} \rVert_1 = \Tr \sqrt{\sqrt{\rho}\,\sigma\sqrt{\rho}} \;\in\; [0,1].
\]
By Uhlmann's theorem~\cite{Uhlmann1976}, $F(\rho, \sigma) = \max \lvert \braket{\phi_\rho}{\phi_\sigma} \rvert$, where the maximum is taken over all purifications $\ket{\phi_\rho}, \ket{\phi_\sigma} \in \cH \otimes \cH_E$ of $\rho$ and $\sigma$ on a sufficiently large purifying space $\cH_E$. We use three standard properties of root fidelity: multiplicativity under tensor products, $F(\rho_1 \otimes \rho_2, \sigma_1 \otimes \sigma_2) = F(\rho_1, \sigma_1) F(\rho_2, \sigma_2)$; additivity over orthogonal classical blocks,
\[
F\!\left(\sum_x p_x \ketbra{x} \otimes \rho_x,\; \sum_x p_x \ketbra{x} \otimes \sigma_x\right) = \sum_x p_x F(\rho_x, \sigma_x)
\]
for any probability distribution $\{p_x\}$ (see~\cite{FuchsVanDeGraaf1999} for the standard equivalence between $F$ and $T_{\mathrm{tr}}$).
(Note that~\cite{WildeGuha2013} defines Uhlmann fidelity as $F(\rho,\sigma)^2$; we account for this convention whenever invoking results from~\cite{WildeGuha2013}.)

\subsection{Binary-Input Classical--Quantum Channels and Polar Combinations}

A \emph{binary-input classical--quantum (c-q) channel} $W : \{0,1\} \to \cD(\cH)$ assigns to each input bit $x \in \{0,1\}$ an output quantum state $\rho_x \in \cD(\cH)$. Under a uniform input $X \sim \Unif(\{0,1\})$, the joint input--output state on $\bC^2 \otimes \cH$ is $\rho_{XB} = \frac{1}{2}\ketbra{0}_X \otimes \rho_0 + \frac{1}{2}\ketbra{1}_X \otimes \rho_1$. The \emph{conditional entropy}, \emph{symmetric capacity}, and \emph{fidelity parameter} of $W$ are defined by
\begin{equation}\label{eq:channel-params}
H(W) := H(X \mid B) = 1 - \QJS(\rho_0, \rho_1), \quad
I(W) := 1 - H(W) = \QJS(\rho_0, \rho_1), \quad
Z(W) := F(\rho_0, \rho_1),
\end{equation}
where $H(X \mid B) := S(\rho_{XB}) - S(\rho_B)$. Since $I(W) = I(X; B)_{\rho} = D(\rho_{XB} \Vert \rho_X \otimes \rho_B)$, the data processing inequality for quantum relative entropy~\cite{Lindblad1975} implies that post-processing the output of $W$ by any CPTP map $\cN$ satisfies $I(\cN \circ W) \le I(W)$, or equivalently $H(\cN \circ W) \ge H(W)$.

Three benchmark families of binary-input channels play a recurring role in our analysis of information combining and duality:
\begin{itemize}
\item A \emph{pure-state channel} $\Pure(f)$ ($f \in [0,1]$) has pure outputs $\rho_x = \ketbra{\phi_x}$ with overlap $\lvert \braket{\phi_0}{\phi_1} \rvert = f$; since $\frac{1}{2}(\ketbra{\phi_0} + \ketbra{\phi_1})$ has eigenvalues $\frac{1 \pm f}{2}$, its conditional entropy is $H(\Pure(f)) = 1 - h_2\!\left(\frac{1-f}{2}\right)$ and its fidelity parameter is $Z(\Pure(f)) = f$.
\item The classical \emph{binary symmetric channel} $\BSC(p)$ ($p \in [0, 1/2]$) has commuting outputs $\rho_0 = \diag(1-p, p)$ and $\rho_1 = \diag(p, 1-p)$, with $H(\BSC(p)) = h_2(p)$ and $Z(\BSC(p)) = 2\sqrt{p(1-p)}$.
\item The classical \emph{binary erasure channel} $\BEC(\epsilon)$ ($\epsilon \in [0,1]$) has commuting outputs $\rho_0 = \diag(1-\epsilon, 0, \epsilon)$ and $\rho_1 = \diag(0, 1-\epsilon, \epsilon)$, with $H(\BEC(\epsilon)) = \epsilon$ and $Z(\BEC(\epsilon)) = \epsilon$.
\end{itemize}

Given two independent binary-input c-q channels $W_1, W_2$, Ar{\i}kan's single-step polar transform maps independent uniform input bits $(U_1, U_2)$ to $(X_1, X_2) = (U_1 \oplus U_2, U_2)$, which are transmitted through $W_1$ and $W_2$ to produce quantum outputs $B_1, B_2$~\cite{Arikan2009,WildeGuha2013}. Decoding $U_1$ from $B_1 B_2$ (treating $U_2$ as uniform noise), and subsequently decoding $U_2$ from $B_1 B_2$ given $U_1$, defines two synthesized binary-input c-q channels.

\begin{definition}[Polar Channel Combinations]\label{def:polar-combinations}
Let $W_1 : x \mapsto \rho_x \in \cD(\cH_1)$ and $W_2 : x \mapsto \sigma_x \in \cD(\cH_2)$ be independent binary-input c-q channels. Their \emph{check-node (minus)} and \emph{variable-node (plus)} combinations are the binary-input c-q channels
\begin{align}
(W_1 \boxtimes W_2)(u_1) &:= \frac{1}{2}\bigl(\rho_{u_1} \otimes \sigma_0 + \rho_{u_1 \oplus 1} \otimes \sigma_1\bigr) \in \cD(\cH_1 \otimes \cH_2), \label{eq:check-node}\\
(W_1 \circledast W_2)(u_2) &:= \frac{1}{2}\ketbra{0}_{U_1} \otimes \rho_{u_2} \otimes \sigma_{u_2} + \frac{1}{2}\ketbra{1}_{U_1} \otimes \rho_{u_2 \oplus 1} \otimes \sigma_{u_2} \in \cD(\bC^2 \otimes \cH_1 \otimes \cH_2). \label{eq:var-node}
\end{align}
For identical copies $W_1 = W_2 = W$, we write $W^- := W \boxtimes W$ and $W^+ := W \circledast W$.
\end{definition}

By conjugating with the SWAP operator on $\cH_1 \otimes \cH_2$, the outputs $(W_1 \boxtimes W_2)(u_1)$ and $(W_2 \boxtimes W_1)(u_1)$ are unitarily equivalent, so $H(W_1 \boxtimes W_2) = H(W_2 \boxtimes W_1)$. Retaining the classical register $U_1$ in~\eqref{eq:var-node} is necessary whenever $W_1$ and $W_2$ are not output-symmetric~\cite{Renes2018,HircheReeb2018}. Because $(W_1 \circledast W_2)(u_2)$ is block-diagonal in the computational basis of $U_1$, its conditional entropy is the average of its two $U_1$-conditioned branches:
\begin{equation}\label{eq:var-node-entropy}
H(W_1 \circledast W_2) = H(U_2 \mid U_1 B_1 B_2) = \frac{1}{2}\bigl[1 - \QJS(\rho_0 \otimes \sigma_0, \rho_1 \otimes \sigma_1)\bigr] + \frac{1}{2}\bigl[1 - \QJS(\rho_1 \otimes \sigma_0, \rho_0 \otimes \sigma_1)\bigr].
\end{equation}
If one omits $U_1$ and evaluates only the $U_1 = 0$ branch $\rho_{u_2} \otimes \sigma_{u_2}$, the two terms in~\eqref{eq:var-node-entropy} can differ on asymmetric channels (such as the classical Z-channel), violating entropy conservation. Under Definition~\ref{def:polar-combinations}, the chain rule for conditional von Neumann entropy and the bijectivity of $(U_1, U_2) \mapsto (U_1 \oplus U_2, U_2)$ yield conservation of entropy for all binary-input c-q channels:
\begin{equation}\label{eq:entropy-conservation}
\begin{aligned}
H(W_1 \boxtimes W_2) + H(W_1 \circledast W_2) &= H(U_1 \mid B_1 B_2) + H(U_2 \mid U_1 B_1 B_2)\\
&= H(U_1 U_2 \mid B_1 B_2) = H(X_1 X_2 \mid B_1 B_2) = H(W_1) + H(W_2).
\end{aligned}
\end{equation}
In particular, $H(W^-) + H(W^+) = 2 H(W)$, so $H(W^-) - H(W) = H(W) - H(W^+)$. Iterating the polar transform $n$ times on $N = 2^n$ independent copies of $W$ produces the $N$ synthesized channels $\{W_n^{(i)}\}_{i=1}^{N}$ associated with Ar{\i}kan's generator matrix $G_N = B_N \left(\begin{smallmatrix} 1 & 0 \\ 1 & 1 \end{smallmatrix}\right)^{\otimes n}$, where $B_N$ is the bit-reversal permutation~\cite{Arikan2009,WildeGuha2013}.

\subsection{Channel Duality}

We also use the channel duality introduced by Renes~\cite{Renes2018} (see also~\cite[Section~IV-A]{HircheReeb2018}). Two binary-input c-q channels $W : \{0,1\} \to \cD(\cH)$ and $W' : \{0,1\} \to \cD(\cH')$ are \emph{reversibly CPTP equivalent} if there exist CPTP maps $\cE : \cB(\cH) \to \cB(\cH')$ and $\cF : \cB(\cH') \to \cB(\cH)$ such that $W' = \cE \circ W$ and $W = \cF \circ W'$ (equivalently, their output pairs agree up to an output isometry and input-independent ancilla states~\cite{Renes2018}). Reversible CPTP equivalence preserves conditional entropy, symmetric capacity, root fidelity, and all polar combinations.

Given a binary-input c-q channel $W : x \mapsto \rho_x \in \cD(\cH)$ with purifications $\ket{\psi_x} \in \cH \otimes \cH_E$ of $\rho_x$, let $V : \bC^2 \to \cH \otimes \cH_E \otimes \bC^2$ be the isometry $V\ket{x} = \ket{\psi_x} \otimes \ket{x}$. Applying $V$ to the Hadamard basis states $\ket{\tilde{x}} := \frac{1}{\sqrt{2}}(\ket{0} + (-1)^x \ket{1})$ and tracing out $\cH$ defines the \emph{dual channel} $W^{\perp} : \{0,1\} \to \cD(\cH_E \otimes \bC^2)$, which is unique up to reversible CPTP equivalence (equivalently, output isometry up to input-independent ancillas).

\begin{proposition}[Renes' Channel Duality~\cite{Renes2018}]\label{prop:duality}
Every binary-input c-q channel $W$ admits a dual binary-input c-q channel $W^{\perp}$ satisfying (with channel equalities in \textup{(D2)--(D3)} understood up to reversible CPTP equivalence, equivalently output isometry up to input-independent ancillas):
\begin{enumerate}
\item[\textup{(D1)}] $H(W) + H(W^{\perp}) = 1$, or equivalently $I(W) + I(W^{\perp}) = 1$;
\item[\textup{(D2)}] $(W_1 \boxtimes W_2)^{\perp} = W_1^{\perp} \circledast W_2^{\perp}$ and $(W_1 \circledast W_2)^{\perp} = W_1^{\perp} \boxtimes W_2^{\perp}$ for all independent binary-input c-q channels $W_1, W_2$;
\item[\textup{(D3)}] $\BSC(p)^{\perp} = \Pure(1 - 2p)$ and $\Pure(1 - 2p)^{\perp} = \BSC(p)$ for all $p \in [0, 1/2]$, while $\BEC(\epsilon)^{\perp} = \BEC(1 - \epsilon)$ for all $\epsilon \in [0, 1]$.
\end{enumerate}
\end{proposition}

Combining~(D1), (D2), and~\eqref{eq:entropy-conservation} gives the duality identity for check-node combining:
\begin{equation}\label{eq:duality-check}
\begin{aligned}
H(W_1 \boxtimes W_2) &= H(W_1) + H(W_2) - H(W_1 \circledast W_2)\\
&= H(W_1) + H(W_2) - \bigl(1 - H((W_1 \circledast W_2)^{\perp})\bigr)\\
&= H(W_1) + H(W_2) - 1 + H(W_1^{\perp} \boxtimes W_2^{\perp}).
\end{aligned}
\end{equation}
Any lower bound on $H(W_1 \boxtimes W_2)$ in terms of $(H_1, H_2)$ therefore produces a dual lower bound under the involution $(H_1, H_2) \mapsto (1 - H_1, 1 - H_2)$. This duality symmetry explains why the conjectured optimal quantum combining bound $g_q(H_1, H_2)$ of~\cite{HircheReeb2018} is the minimum of the classical bound $g_c(H_1, H_2)$ (attained by binary symmetric channels) and its dual $H_1 + H_2 - 1 + g_c(1 - H_1, 1 - H_2)$ (attained by pure-state channels).

\section{Quantum Jeffreys--Jensen--Shannon Inequality and Dephasing Contraction}\label{sec:qjjs-sdpi}

In this section we prove Theorem~\ref{thm:dephasing-sdpi}. The central analytic ingredient is a comparison inequality (Proposition~\ref{prop:qjjs}) between the symmetrized quantum relative entropy $\frac{1}{4}(D(\rho\Vert\sigma) + D(\sigma\Vert\rho))$ and the quantum Jensen--Shannon divergence $\QJS(\rho, \sigma)$. In the classical setting, inequalities comparing two $f$-divergences are obtained by verifying a scalar inequality for their generating functions and integrating against the likelihood-ratio distribution~\cite{SasonVerdu2016}. When $\rho$ and $\sigma$ do not commute, no single likelihood-ratio distribution exists directly, and operator Jensen arguments lose tightness. We circumvent this obstruction by using Frenkel's integral representation~\cite{Frenkel2023} in its hockey-stick form~\cite{HircheTomamichel2024} to show that $D(\rho\Vert\sigma)$, $D(\sigma\Vert\rho)$, and $\QJS(\rho, \sigma)$ are all determined by a single convex scalar function $s \mapsto E_s(\rho\Vert\sigma)$, and therefore admit simultaneous representations as classical $f$-divergences against one compactly supported probability measure on $(0,\infty)$.

\subsection{Scalar Convexity and Simultaneous Hockey-Stick Representation}

Define $\phi : [0,1] \to [0,\infty]$ on $[0,1)$ parametrically by
\begin{equation}\label{eq:phi-def}
\phi(c) := \frac{x}{2}\log_2\!\left(\frac{1+x}{1-x}\right), \qquad
c = 1 - h_2\!\left(\frac{1-x}{2}\right), \qquad x \in [0,1),
\end{equation}
with $\phi(1) := +\infty$. Geometrically, $c(x) = \QJS(\rho_x, \sigma_x)$ and $2\phi(c(x)) = D(\rho_x \Vert \sigma_x) = D(\sigma_x \Vert \rho_x)$ for the symmetric binary pair $\rho_x = \diag(\frac{1+x}{2}, \frac{1-x}{2})$, $\sigma_x = \diag(\frac{1-x}{2}, \frac{1+x}{2})$. We first establish the strict convexity of $\phi$ and express its supporting tangent inequalities in perspective-symmetric form.

\begin{lemma}[Perspective-Symmetric Scalar Inequality]\label{lem:scalar-convexity}
The function $\phi$ defined in~\eqref{eq:phi-def} is continuous, strictly increasing, and strictly convex on $[0,1)$, with derivative
\begin{equation}\label{eq:phi-deriv}
\phi'(c(x)) = 1 + \frac{2x}{(1 - x^2)\ln\!\left(\frac{1+x}{1-x}\right)}, \qquad x \in (0,1).
\end{equation}
Moreover, for every $x \in (0,1)$, setting $\lambda(x) := \phi'(c(x))$ and $\beta(x) := \lambda(x)c(x) - \phi(c(x))$, the function
\begin{equation}\label{eq:gx-def}
g_x(t) := \frac{1}{4}(t-1)\log_2 t - \lambda(x)\left(\frac{t}{2}\log_2\!\left(\frac{2t}{1+t}\right) + \frac{1}{2}\log_2\!\left(\frac{2}{1+t}\right)\right) + \beta(x)\frac{1+t}{2}
\end{equation}
satisfies $g_x(t) \ge 0$ for all $t > 0$, with equality if and only if $t \in \bigl\{\frac{1+x}{1-x}, \frac{1-x}{1+x}\bigr\}$.
\end{lemma}

\begin{proof}
For $x \in (0,1)$, rewrite $c(x)$ in symmetric form as
\[
c(x) = 1 - h_2\!\left(\frac{1-x}{2}\right) = \frac{1+x}{2}\log_2(1+x) + \frac{1-x}{2}\log_2(1-x).
\]
Differentiating $c(x)$ and $\phi(c(x)) = \frac{x}{2}\log_2(\frac{1+x}{1-x})$ with respect to $x \in (0,1)$ yields
\[
\frac{dc}{dx} = \frac{1}{2}\log_2\!\left(\frac{1+x}{1-x}\right) > 0, \qquad
\frac{d\phi}{dx} = \frac{1}{2}\log_2\!\left(\frac{1+x}{1-x}\right) + \frac{x}{(\ln 2)(1 - x^2)}.
\]
Thus $x \mapsto c(x)$ is a smooth, strictly increasing bijection from $[0,1)$ onto $[0,1)$, and taking the ratio $(d\phi/dx)/(dc/dx)$ gives~\eqref{eq:phi-deriv}. To verify strict convexity of $\phi$, substitute $x = \tanh u$ ($u > 0$), so that $\frac{1+x}{1-x} = e^{2u}$, $\ln(\frac{1+x}{1-x}) = 2u$, and $\frac{2x}{1-x^2} = 2\sinh u \cosh u = \sinh(2u)$. Then
\[
\phi'\bigl(c(\tanh u)\bigr) = 1 + \frac{\sinh(2u)}{2u} = 1 + \sum_{k=0}^{\infty} \frac{(2u)^{2k}}{(2k+1)!},
\]
which is strictly increasing on $(0,\infty)$ because every term in the power series with $k \ge 1$ is strictly increasing. Since $u \mapsto c(\tanh u)$ is a strictly increasing bijection from $(0,\infty)$ onto $(0,1)$, $\phi'$ is strictly increasing on $(0,1)$, and hence $\phi$ is continuous, strictly increasing, and strictly convex on $[0,1)$.

For any $t > 0$, introduce the Cayley variable $y := \frac{t-1}{t+1} \in (-1, 1)$, so that $t = \frac{1+y}{1-y}$, $\frac{1+t}{2} = \frac{1}{1-y}$, $\frac{2t}{1+t} = 1+y$, and $\frac{2}{1+t} = 1-y$. Dividing~\eqref{eq:gx-def} by $\frac{1+t}{2} > 0$ gives
\begin{align*}
\frac{g_x(t)}{(1+t)/2} &= \frac{1}{2}\cdot\frac{t-1}{t+1}\log_2 t - \lambda(x)\left(\frac{t}{1+t}\log_2\!\left(\frac{2t}{1+t}\right) + \frac{1}{1+t}\log_2\!\left(\frac{2}{1+t}\right)\right) + \beta(x)\\
&= \frac{y}{2}\log_2\!\left(\frac{1+y}{1-y}\right) - \lambda(x)\left(\frac{1+y}{2}\log_2(1+y) + \frac{1-y}{2}\log_2(1-y)\right) + \beta(x).
\end{align*}
Both $y \mapsto \frac{y}{2}\log_2(\frac{1+y}{1-y})$ and $y \mapsto \frac{1+y}{2}\log_2(1+y) + \frac{1-y}{2}\log_2(1-y)$ are even functions on $(-1,1)$ vanishing at $y=0$, so they coincide with $\phi(c(\lvert y \rvert))$ and $c(\lvert y \rvert)$, respectively. Substituting $\beta(x) = \lambda(x)c(x) - \phi(c(x))$ and $\lambda(x) = \phi'(c(x))$ therefore yields
\[
\frac{g_x(t)}{(1+t)/2} = \phi\bigl(c(\lvert y \rvert)\bigr) - \phi\bigl(c(x)\bigr) - \phi'\bigl(c(x)\bigr)\bigl(c(\lvert y \rvert) - c(x)\bigr) \;\ge\; 0
\]
by the strict convexity of $\phi$ on $[0,1)$. Equality holds if and only if $c(\lvert y \rvert) = c(x)$, i.e., $\lvert y \rvert = x$, which is equivalent to $t \in \bigl\{\frac{1+x}{1-x}, \frac{1-x}{1+x}\bigr\}$.
\end{proof}

We next construct the compactly supported Borel probability measure $\mu$ on $(0,\infty)$ that simultaneously represents $D(\rho\Vert\sigma)$, $D(\sigma\Vert\rho)$, and $\QJS(\rho, \sigma)$. The key observation is that the hockey-stick divergence of $\rho$ relative to the midpoint $m = \frac{\rho+\sigma}{2}$ reduces algebraically to a rescaled hockey-stick divergence between $\rho$ and $\sigma$ itself.

\begin{lemma}[Simultaneous Hockey-Stick Representation]\label{lem:hockey-stick-measure}
Let $\rho, \sigma \in \cD(\cH)$ satisfy $\supp(\rho) = \supp(\sigma)$. Then there exists a compactly supported Borel probability measure $\mu$ on $[a, b] \subset (0, \infty)$ with $\int_a^b t\, d\mu(t) = 1$ such that $E_s(\rho \Vert \sigma) = \int_a^b (t - s)_+\, d\mu(t)$ for all $s \ge 0$, and simultaneously
\begin{align}
D(\rho \Vert \sigma) &= \int_a^b \left(t\log_2 t - \frac{t-1}{\ln 2}\right) d\mu(t), \label{eq:D-rho-sigma-rep}\\
D(\sigma \Vert \rho) &= \int_a^b \left(-\log_2 t + \frac{t-1}{\ln 2}\right) d\mu(t), \label{eq:D-sigma-rho-rep}\\
\QJS(\rho, \sigma) &= \int_a^b \left(\frac{t}{2}\log_2\!\left(\frac{2t}{1+t}\right) + \frac{1}{2}\log_2\!\left(\frac{2}{1+t}\right)\right) d\mu(t). \label{eq:qjs-rep}
\end{align}
\end{lemma}

\begin{proof}
Restricting both states to their common support $\cH_0 := \supp(\rho) = \supp(\sigma)$, we may assume that $\rho$ and $\sigma$ are strictly positive definite on $\cH_0$. Let $a := \lambda_{\min}(\sigma^{-1/2}\rho\sigma^{-1/2}) > 0$ and $b := \lambda_{\max}(\sigma^{-1/2}\rho\sigma^{-1/2}) < \infty$, so that $a\sigma \preceq \rho \preceq b\sigma$ on $\cH_0$. By~\eqref{eq:hockey-stick-def},
\[
f(s) := E_s(\rho \Vert \sigma) = \max_{0 \preceq \Pi \preceq I} \Tr\!\bigl[\Pi(\rho - s\sigma)\bigr], \qquad s \ge 0,
\]
is the pointwise supremum of affine functions of $s$, and is therefore convex and continuous on $[0,\infty)$. For $s \in [0,a]$, the operator $\rho - s\sigma$ is positive semidefinite, so the maximum is attained at $\Pi = I$ and $f(s) = 1 - s$. For $s \ge b$, $\rho - s\sigma \preceq 0$, so the maximum is attained at $\Pi = 0$ and $f(s) = 0$.

Consequently, the right derivative $f'_+(s)$ is nondecreasing and right-continuous on $[0,\infty)$, equals $-1$ on $[0,a)$, and equals $0$ on $[b,\infty)$. Its Lebesgue--Stieltjes measure $d\mu(t) := d(f'_+(t))$ is therefore a nonnegative Borel measure supported on $[a,b]$ with total mass $\mu([a,b]) = f'_+(b) - f'_+(0) = 1$. Since $-f'_+(u) = f'_+(b) - f'_+(u) = \int_{(u,b]} d\mu(t)$ for every $u \ge 0$ and $f(b) = 0$, the fundamental theorem of calculus for convex functions and Tonelli's theorem give
\[
f(s) = -\int_s^b f'_+(u)\, du = \int_s^b \left(\int_{(u,b]} d\mu(t)\right) du = \int_{[a,b]} \left(\int_s^b \mathbf{1}_{\{u < t\}}\, du\right) d\mu(t) = \int_a^b (t - s)_+\, d\mu(t)
\]
for all $s \in [0, b]$, regardless of whether $\mu$ has atoms at $a$ or $b$; for $s > b$, both $f(s)$ and $\int_a^b (t - s)_+\, d\mu(t)$ vanish identically. Evaluating at $s = 0$ yields $\int_a^b t\, d\mu(t) = f(0) = 1$, and hence $\int_a^b (t - 1)\, d\mu(t) = 0$.

By Frenkel's integral representation~\cite[Theorem~6]{Frenkel2023} (converted to bits), for any $\rho, \sigma \in \cD(\cH)$ with $\supp(\rho) \subseteq \supp(\sigma)$,
\begin{equation}\label{eq:frenkel-integral}
D(\rho \Vert \sigma) = \frac{1}{\ln 2} \int_{-\infty}^{\infty} \frac{\Tr\!\bigl(((1-u)\rho + u\sigma)_-\bigr)}{\lvert u \rvert (u - 1)^2}\, du.
\end{equation}
On $[0,1]$, the convex combination $(1-u)\rho + u\sigma$ is positive semidefinite, so the integrand vanishes. We split the remaining domain into $(-\infty, 0)$ and $(1,\infty)$:
\begin{itemize}
\item On $(-\infty, 0)$, set $s = \frac{-u}{1-u} \in (0,1)$, so $1-u = \frac{1}{1-s}$ and $du = -\frac{ds}{(1-s)^2}$. Since $(1-u)\rho + u\sigma = \frac{1}{1-s}(\rho - s\sigma)$ and $\Tr[(\rho - s\sigma)_-] = \Tr[(\rho - s\sigma)_+] - \Tr[\rho - s\sigma] = E_s(\rho\Vert\sigma) - (1-s)$, the integrand becomes $\frac{E_s(\rho\Vert\sigma) - (1-s)}{s}\,ds$.
\item On $(1,\infty)$, set $s = \frac{u}{u-1} \in (1,\infty)$, so $u-1 = \frac{1}{s-1}$ and $du = -\frac{ds}{(s-1)^2}$. Since $(1-u)\rho + u\sigma = -\frac{1}{s-1}(\rho - s\sigma)$ gives $\Tr[((1-u)\rho + u\sigma)_-] = \frac{1}{s-1} E_s(\rho\Vert\sigma)$, the integrand becomes $\frac{E_s(\rho\Vert\sigma)}{s}\,ds$.
\end{itemize}
Combining both intervals transforms~\eqref{eq:frenkel-integral} into the hockey-stick representation~\cite{HircheTomamichel2024} (see also~\cite{BertaLamiTomamichel2024}):
\begin{equation}\label{eq:hockey-stick-integral}
D(\rho \Vert \sigma) = \frac{1}{\ln 2} \int_0^{\infty} \frac{E_s(\rho \Vert \sigma) - (1 - s)_+}{s}\, ds.
\end{equation}

For $s \in (0,1]$ and $t > 0$, the identity $(t - s)_+ - (1 - s) = (s - t)_+ + (t - 1)$ holds pointwise. Integrating against $d\mu(t)$ and using $\int_a^b (t - 1)\, d\mu(t) = 0$ yields
\[
E_s(\rho \Vert \sigma) - (1 - s)_+ =
\begin{cases}
\int_a^b (s - t)_+\, d\mu(t), & s \in (0, 1],\\
\int_a^b (t - s)_+\, d\mu(t), & s > 1,
\end{cases}
\]
with nonnegative integrands throughout. Substituting into~\eqref{eq:hockey-stick-integral} and interchanging the integrals by Tonelli's theorem gives
\[
\int_0^1 \frac{(s - t)_+}{s}\, ds + \int_1^{\infty} \frac{(t - s)_+}{s}\, ds =
\begin{cases}
\int_t^1 \bigl(1 - \frac{t}{s}\bigr)\, ds = (1 - t) + t\ln t, & 0 < t \le 1,\\
\int_1^t \bigl(\frac{t}{s} - 1\bigr)\, ds = t\ln t - (t - 1), & t > 1.
\end{cases}
\]
In both cases the inner integral equals $t\ln t - t + 1$, and dividing by $\ln 2$ proves~\eqref{eq:D-rho-sigma-rep}.

Next, for every $s > 0$, the trace identity $\Tr[(\sigma - s\rho)_+] = \Tr[(s\rho - \sigma)_+] + \Tr[\sigma - s\rho]$ gives $E_s(\sigma \Vert \rho) = s E_{1/s}(\rho \Vert \sigma) + 1 - s$. Substituting $E_{1/s}(\rho \Vert \sigma) = \int_a^b (t - 1/s)_+\, d\mu(t)$ and applying $(st - 1)_+ + (1 - st) = (1 - st)_+$ together with $\int_a^b s(t - 1)\, d\mu(t) = 0$ yields $E_s(\sigma \Vert \rho) = \int_a^b (1 - st)_+\, d\mu(t)$. On $s \in (0,1]$, writing $(1 - st)_+ - (1 - s) = (st - 1)_+ + s(1 - t)$ and integrating against $d\mu(t)$ yields $E_s(\sigma \Vert \rho) - (1 - s) = \int_a^b (st - 1)_+\, d\mu(t)$. Substituting into~\eqref{eq:hockey-stick-integral}, applying Tonelli's theorem, and changing variables to $u = st$ gives
\[
\int_0^1 \frac{(st - 1)_+}{s}\, ds + \int_1^{\infty} \frac{(1 - st)_+}{s}\, ds = \int_0^t \frac{(u - 1)_+}{u}\, du + \int_t^{\infty} \frac{(1 - u)_+}{u}\, du = -\ln t + t - 1
\]
for every $t > 0$, establishing~\eqref{eq:D-sigma-rho-rep}.

Finally, let $m := \frac{\rho + \sigma}{2}$. For $s > 0$, we have $\rho - sm = (1 - \frac{s}{2})\rho - \frac{s}{2}\sigma$. If $s \ge 2$, then $\rho - sm \preceq 0$ and hence $E_s(\rho \Vert m) = 0$. If $s \in (0,2)$, then $1 - \frac{s}{2} > 0$, so positive homogeneity of the positive part gives
\begin{equation}\label{eq:midpoint-hockey}
E_s(\rho \Vert m) = \left(1 - \frac{s}{2}\right) E_{\frac{s}{2-s}}(\rho \Vert \sigma) = \int_a^b \left(t - s\frac{1+t}{2}\right)_+ d\mu(t).
\end{equation}
By the same calculation for $\sigma - sm = (1 - \frac{s}{2})\sigma - \frac{s}{2}\rho$ using $E_\gamma(\sigma \Vert \rho) = \int_a^b (1 - \gamma t)_+\, d\mu(t)$, we obtain $E_s(\sigma \Vert m) = \int_a^b (1 - s\frac{1+t}{2})_+\, d\mu(t)$ for $s \in (0,2)$ and $0$ for $s \ge 2$. On $s \in (0,1]$,
\begin{align*}
\left(t - s\frac{1+t}{2}\right)_+ - (1 - s) &= \left(s\frac{1+t}{2} - t\right)_+ + (t - 1)\left(1 - \frac{s}{2}\right),\\
\left(1 - s\frac{1+t}{2}\right)_+ - (1 - s) &= \left(s\frac{1+t}{2} - 1\right)_+ + (1 - t)\frac{s}{2},
\end{align*}
whose linear remainders vanish upon integration against $d\mu(t)$ since $\int_a^b (t - 1)\, d\mu(t) = 0$. Substituting into~\eqref{eq:hockey-stick-integral}, applying Tonelli's theorem, and setting $r := \frac{2t}{1+t} \in (0,2)$ gives
\begin{align*}
\int_0^1 \frac{(s\frac{1+t}{2} - t)_+}{s}\, ds + \int_1^2 \frac{(t - s\frac{1+t}{2})_+}{s}\, ds
&= \frac{1+t}{2}\left(\int_0^1 \frac{(s - r)_+}{s}\, ds + \int_1^2 \frac{(r - s)_+}{s}\, ds\right)\\
&= t\ln\!\left(\frac{2t}{1+t}\right) + \frac{1 - t}{2},
\end{align*}
and symmetrically with $t$ and $1$ exchanged,
\[
\int_0^1 \frac{(s\frac{1+t}{2} - 1)_+}{s}\, ds + \int_1^2 \frac{(1 - s\frac{1+t}{2})_+}{s}\, ds
= \ln\!\left(\frac{2}{1+t}\right) + \frac{t - 1}{2}.
\]
Averaging $D(\rho \Vert m)$ and $D(\sigma \Vert m)$ cancels the $\pm\frac{1-t}{2}$ terms and yields~\eqref{eq:qjs-rep}.
\end{proof}

\subsection{Quantum Jeffreys--Jensen--Shannon Inequality and Proof of Theorem~\ref{thm:dephasing-sdpi}}

Integrating the scalar inequality of Lemma~\ref{lem:scalar-convexity} against the simultaneous probability measure $\mu$ of Lemma~\ref{lem:hockey-stick-measure} yields the desired comparison inequality between the symmetrized quantum relative entropy (quantum Jeffreys divergence) and the quantum Jensen--Shannon divergence.

\begin{proposition}[Quantum Jeffreys--Jensen--Shannon Inequality]\label{prop:qjjs}
For every finite-dimensional Hilbert space $\cH$ and all density operators $\rho, \sigma \in \cD(\cH)$,
\begin{equation}\label{eq:qjjs}
\frac{1}{4}\bigl(D(\rho \Vert \sigma) + D(\sigma \Vert \rho)\bigr) \;\ge\; \phi\bigl(\QJS(\rho, \sigma)\bigr),
\end{equation}
where $\phi : [0,1] \to [0,\infty]$ is defined in~\eqref{eq:phi-def}. Equality holds for every symmetric classical binary pair $\rho = \diag(\frac{1+x}{2}, \frac{1-x}{2})$, $\sigma = \diag(\frac{1-x}{2}, \frac{1+x}{2})$ with $x \in [0,1]$.
\end{proposition}

\begin{proof}
If $\supp(\rho) \ne \supp(\sigma)$, then $D(\rho \Vert \sigma) + D(\sigma \Vert \rho) = +\infty$ and~\eqref{eq:qjjs} holds (including when $\rho \perp \sigma$, in which case $\QJS(\rho, \sigma) = 1$ and $\phi(1) = +\infty$). If $\rho = \sigma$, both sides vanish.

Assume therefore that $\supp(\rho) = \supp(\sigma)$ and $\rho \ne \sigma$, so that $c := \QJS(\rho, \sigma) \in (0,1)$. Let $x \in (0,1)$ be the unique solution to $c(x) = c$, and let $\mu$ be the compactly supported Borel probability measure on $[a,b] \subset (0,\infty)$ given by Lemma~\ref{lem:hockey-stick-measure}. Adding~\eqref{eq:D-rho-sigma-rep} and~\eqref{eq:D-sigma-rho-rep} cancels the linear terms $\pm\frac{t-1}{\ln 2}$ and gives
\[
\frac{1}{4}\bigl(D(\rho \Vert \sigma) + D(\sigma \Vert \rho)\bigr) = \int_a^b \frac{1}{4}(t - 1)\log_2 t\, d\mu(t).
\]
Since $\mu$ is a probability measure with $\int_a^b t\, d\mu(t) = 1$, we have $\int_a^b \frac{1+t}{2}\, d\mu(t) = 1$. Integrating the pointwise inequality $g_x(t) \ge 0$ of Lemma~\ref{lem:scalar-convexity} against $\mu$ and substituting~\eqref{eq:qjs-rep} yields
\[
\frac{1}{4}\bigl(D(\rho \Vert \sigma) + D(\sigma \Vert \rho)\bigr) - \lambda(x)\QJS(\rho, \sigma) + \beta(x) = \int_a^b g_x(t)\, d\mu(t) \;\ge\; 0.
\]
By definition of $\beta(x) = \lambda(x)c(x) - \phi(c(x))$, we have $\lambda(x)\QJS(\rho, \sigma) - \beta(x) = \phi(c(x)) = \phi(\QJS(\rho, \sigma))$, which establishes~\eqref{eq:qjjs}. For the symmetric classical binary pair $\rho = \diag(\frac{1+x}{2}, \frac{1-x}{2})$ and $\sigma = \diag(\frac{1-x}{2}, \frac{1+x}{2})$ with $x \in [0,1)$, direct computation gives $\QJS(\rho, \sigma) = c(x)$ and $D(\rho \Vert \sigma) = D(\sigma \Vert \rho) = 2\phi(c(x))$, attaining equality in~\eqref{eq:qjjs}; at $x = 1$, both sides equal $+\infty$.
\end{proof}

\begin{remark}
Proposition~\ref{prop:qjjs} is a noncommutative counterpart of classical $f$-divergence range bounds between Jeffreys divergence and Jensen--Shannon divergence (see~\cite{SasonVerdu2016}), and complements lower bounds on the concavity of von Neumann entropy in terms of relative entropy and trace distance~\cite{KimRuskai2014,MullerHermesFrancaWolf2016}. Note that $\phi(c)/c \to 2$ as $c \to 0^+$, while $\phi(c)$ diverges as $c \to 1^-$; this boundary divergence is what drives the nonlinear contraction rate in Theorem~\ref{thm:dephasing-sdpi}.
\end{remark}

We now prove Theorem~\ref{thm:dephasing-sdpi} by applying Proposition~\ref{prop:qjjs} to the entropy production of the continuous-time qubit dephasing semigroup, in analogy with modified logarithmic Sobolev arguments for quantum Markov semigroups~\cite{KastoryanoTemme2013,MullerHermesFrancaWolf2016}.

\begin{proof}[Proof of Theorem~\ref{thm:dephasing-sdpi}]
Consider the continuous-time qubit dephasing semigroup $\cZ_{p(t)} = e^{t\cL}$ on $\cB(\bC^2)$ with Lindblad generator $\cL(X) = \frac{1}{2}(ZXZ - X)$. Since $Z^2 = I$, we have $\cL^2 = -\cL$, and expanding the exponential series gives $e^{t\cL}(X) = (1 - p(t))X + p(t)ZXZ$ with $p(t) = \frac{1 - e^{-t}}{2} \in [0, 1/2)$ for $t \ge 0$. At $p = 1/2$, $(\cZ_{1/2} \otimes \id)(\Psi_{AB}) = \Psi_\infty$, so $C(1/2) = 0 = 1 - h_2(1/2)$ and~\eqref{eq:dephasing-sdpi} holds with equality. It therefore suffices to consider $p \in [0, 1/2)$.

For any state $\Psi_{AB} \in \cD(\bC^2 \otimes \cH_B)$ and $p \in [0, 1/2]$, set $\Psi_p := (\cZ_p \otimes \id)(\Psi_{AB})$, $\widetilde{\Psi}_p := (Z \otimes I)\Psi_p(Z \otimes I)$, and $\Psi_\infty := (\cZ_{1/2} \otimes \id)(\Psi_{AB}) = \frac{1}{2}(\Psi_p + \widetilde{\Psi}_p)$. Because $\widetilde{\Psi}_p \succeq 0$, we always have $\supp(\Psi_p) \subseteq \supp(\Psi_\infty)$. Moreover, $\cP := \cZ_{1/2} \otimes \id$ is a self-adjoint conditional expectation onto the fixed-point subalgebra of $\cZ_p \otimes \id$ satisfying $\cP(\Psi_p) = \Psi_\infty$ and $\cP(\log_2 \Psi_\infty) = \log_2 \Psi_\infty$ (with $\log_2 \Psi_\infty$ defined on $\supp(\Psi_\infty)$). Hence $\Tr[\Psi_p \log_2 \Psi_\infty] = \Tr[\cP(\Psi_p)\log_2 \Psi_\infty] = -S(\Psi_\infty)$, and therefore
\[
C(p) = D(\Psi_p \Vert \Psi_\infty) = S(\Psi_\infty) - S(\Psi_p).
\]
In particular, $C(p)$ depends continuously on $\Psi_{AB} \in \cD(\bC^2 \otimes \cH_B)$. We may therefore assume first that $\Psi_{AB} \succ 0$ is strictly positive definite; the general case follows by applying~\eqref{eq:dephasing-sdpi} to $\Psi^{(\epsilon)} := (1 - \epsilon)\Psi_{AB} + \epsilon \frac{I}{2\dim \cH_B}$ for $\epsilon \in (0,1)$ and taking $\epsilon \downarrow 0$, using the continuity of von Neumann entropy on $\cD(\bC^2 \otimes \cH_B)$ and of $h_2, h_2^{-1}$ on $[0,1]$.

For $t \ge 0$, write $\Psi_t := (\cZ_{p(t)} \otimes \id)(\Psi_{AB}) \succ 0$ and $\widetilde{\Psi}_t := (Z \otimes I)\Psi_t(Z \otimes I) \succ 0$. Since $\widetilde{\Psi}_t$ is unitarily equivalent to $\Psi_t$, we have $S(\widetilde{\Psi}_t) = S(\Psi_t)$ and hence
\[
\QJS(\Psi_t, \widetilde{\Psi}_t) = S(\Psi_\infty) - S(\Psi_t) = C(p(t)), \qquad
D(\widetilde{\Psi}_t \Vert \Psi_t) = D(\Psi_t \Vert \widetilde{\Psi}_t).
\]
Because $t \mapsto \Psi_t$ is smooth on $[0,\infty)$ with spectrum bounded strictly away from~$0$, differentiating $S(\Psi_t) = -\Tr[\Psi_t \log_2 \Psi_t]$ using $\frac{d}{dt}\Tr[\Psi_t \log_2 \Psi_t] = \Tr[\dot{\Psi}_t \log_2 \Psi_t] + \frac{1}{\ln 2}\Tr[\dot{\Psi}_t]$ and $\Tr[(\cL \otimes \id)(\Psi_t)] = 0$ gives
\[
-\frac{d}{dt} C(p(t)) = \frac{d}{dt} S(\Psi_t) = -\Tr\!\bigl[(\cL \otimes \id)(\Psi_t)\log_2 \Psi_t\bigr] = \frac{1}{2}\Tr\!\bigl[(\Psi_t - \widetilde{\Psi}_t)\log_2 \Psi_t\bigr].
\]
Conjugating inside the trace by the unitary involution $Z \otimes I$ maps $\Psi_t - \widetilde{\Psi}_t$ to $-(\Psi_t - \widetilde{\Psi}_t)$ and $\log_2 \Psi_t$ to $\log_2 \widetilde{\Psi}_t$, so $\Tr[(\Psi_t - \widetilde{\Psi}_t)\log_2 \Psi_t] = -\Tr[(\Psi_t - \widetilde{\Psi}_t)\log_2 \widetilde{\Psi}_t]$. Averaging the two expressions and invoking Proposition~\ref{prop:qjjs} gives
\[
-\frac{d}{dt} C(p(t)) = \frac{1}{4}\Tr\!\bigl[(\Psi_t - \widetilde{\Psi}_t)(\log_2 \Psi_t - \log_2 \widetilde{\Psi}_t)\bigr] = \frac{1}{4}\bigl(D(\Psi_t \Vert \widetilde{\Psi}_t) + D(\widetilde{\Psi}_t \Vert \Psi_t)\bigr) \;\ge\; \phi\bigl(C(p(t))\bigr).
\]

If $C(0) = 0$, then $\Psi_0 = \Psi_\infty$, so $C(p) = 0$ for all $p \in [0, 1/2]$ and~\eqref{eq:dephasing-sdpi} holds directly. Assume therefore that $C(0) > 0$, so $\Psi_0 \ne \widetilde{\Psi}_0$. Since $\Psi_t - \widetilde{\Psi}_t = e^{-t}(\Psi_0 - \widetilde{\Psi}_0) \ne 0$ for every finite $t \ge 0$, we have $\Psi_t \ne \Psi_\infty$ and hence $C(p(t)) = D(\Psi_t \Vert \Psi_\infty) > 0$. On the other hand, $\Psi_t \succ 0$ implies $\supp(\Psi_t) = \supp(\widetilde{\Psi}_t)$, so $C(p(t)) = \QJS(\Psi_t, \widetilde{\Psi}_t) < 1$.

Consequently, $x(t) := 1 - 2h_2^{-1}(1 - C(p(t))) \in (0,1)$ is smooth on $[0,\infty)$ and satisfies $C(p(t)) = 1 - h_2(\frac{1 - x(t)}{2})$. By the chain rule and the derivative formula in Lemma~\ref{lem:scalar-convexity},
\[
\frac{dC}{dt} = \frac{dC}{dx}\frac{dx}{dt} = \frac{1}{2}\log_2\!\left(\frac{1 + x(t)}{1 - x(t)}\right)\frac{dx}{dt} = \frac{\phi(C(p(t)))}{x(t)}\frac{dx}{dt}.
\]
Dividing $-\frac{dC}{dt} \ge \phi(C(p(t))) > 0$ by the strictly positive quantity $\phi(C(p(t)))/x(t)$ reduces the differential inequality to
\[
\frac{dx}{dt} \;\le\; -x(t), \qquad t \ge 0.
\]
Hence $\frac{d}{dt}(e^t x(t)) \le 0$, so $x(t) \le e^{-t} x(0) = (1 - 2p(t)) x(0)$ for all $t \ge 0$. Substituting $x(t) = 1 - 2h_2^{-1}(1 - C(p(t)))$ and applying the elementary identity $(1 - 2p)(1 - 2a) = 1 - 2(a * p)$ yields
\[
1 - 2h_2^{-1}\bigl(1 - C(p)\bigr) \;\le\; (1 - 2p)\bigl(1 - 2h_2^{-1}(1 - C(0))\bigr) = 1 - 2\bigl(h_2^{-1}(1 - C(0)) * p\bigr),
\]
so $h_2^{-1}(1 - C(p)) \ge h_2^{-1}(1 - C(0)) * p$. Applying the increasing function $h_2$ on $[0, 1/2]$ yields~\eqref{eq:dephasing-sdpi}.

Finally, for any equatorial single-qubit state $\Psi_A = \frac{1}{2}(I + xX)$ ($x \in [0,1]$), we have $\cZ_p(\Psi_A) = \frac{1}{2}(I + (1 - 2p)xX)$ and $\cZ_{1/2}(\Psi_A) = \frac{1}{2}I$. The eigenvalues of $\cZ_p(\Psi_A)$ are $\frac{1 \pm (1 - 2p)x}{2}$, so $C(p) = 1 - h_2(\frac{1 - (1 - 2p)x}{2}) = 1 - h_2(\frac{1-x}{2} * p)$, attaining equality in~\eqref{eq:dephasing-sdpi} for every $p \in [0, 1/2]$.
\end{proof}

\section{Quantum Mrs.\ Gerber Bounds}\label{sec:qmgl}

We now apply Theorem~\ref{thm:dephasing-sdpi} to binary quantum information combining and prove Theorem~\ref{thm:qmgl}. The bridge between polar channel combining and local qubit dephasing is a Gram-matrix factorization that identifies the variable-node combination $W_1 \circledast W_2$ with a pure-state channel $W_2$ as a partially dephased bipartite state on $\bC^2 \otimes \cH_1$.

\subsection{Gram-Matrix Dephasing Reduction and Hybrid Pairs}

To relate binary channel combining to qubit dephasing, we factor the mixture $\frac{1}{2}(\rho_0 \otimes \ketbra{\phi_0} + \rho_1 \otimes \ketbra{\phi_1})$ as $A^{\dagger}A$ and pass to the operator $AA^{\dagger}$ on $\bC^2 \otimes \cH_1$, which has the same nonzero spectrum.

\begin{lemma}[Gram-Matrix Dephasing Reduction]\label{lem:gram-dephasing}
Let $\rho_0, \rho_1 \in \cD(\cH_1)$ and let $\ket{\phi_0}, \ket{\phi_1} \in \cH_2$ be unit vectors with overlap $f_2 := \lvert \braket{\phi_0}{\phi_1} \rvert \in [0,1]$. Set $p := \frac{1 - f_2}{2} \in [0, 1/2]$ and define the operator $\Psi_1$ on $\bC^2 \otimes \cH_1$ by
\begin{equation}\label{eq:gram-state}
\Psi_1 := \frac{1}{2}
\begin{pmatrix}
\rho_0 & \sqrt{\rho_0}\sqrt{\rho_1}\\
\sqrt{\rho_1}\sqrt{\rho_0} & \rho_1
\end{pmatrix}.
\end{equation}
Then $\Psi_1 \in \cD(\bC^2 \otimes \cH_1)$, the nonzero eigenvalues (with multiplicities) of $\frac{1}{2}(\rho_0 \otimes \ketbra{\phi_0} + \rho_1 \otimes \ketbra{\phi_1})$ coincide with those of $(\cZ_p \otimes \id)(\Psi_1)$, and
\[
S(\Psi_1) = S\!\left(\frac{\rho_0 + \rho_1}{2}\right), \qquad
S\bigl((\cZ_{1/2} \otimes \id)(\Psi_1)\bigr) = 1 + \frac{1}{2}\bigl(S(\rho_0) + S(\rho_1)\bigr),
\]
so that $C(0) = D(\Psi_1 \Vert (\cZ_{1/2} \otimes \id)(\Psi_1)) = H(W_1)$ for the channel $W_1 : x \mapsto \rho_x$.
\end{lemma}

\begin{proof}
Adjusting the global phase of $\ket{\phi_1}$ so that $\braket{\phi_0}{\phi_1} = f_2 \ge 0$, define the linear operator $A : \cH_1 \otimes \cH_2 \to \bC^2 \otimes \cH_1$ by
\[
A := \frac{1}{\sqrt{2}}
\begin{pmatrix}
\sqrt{\rho_0} \otimes \bra{\phi_0}\\
\sqrt{\rho_1} \otimes \bra{\phi_1}
\end{pmatrix}.
\]
Direct block multiplication gives $A^{\dagger}A = \frac{1}{2}(\rho_0 \otimes \ketbra{\phi_0} + \rho_1 \otimes \ketbra{\phi_1})$, whereas
\[
AA^{\dagger} = \frac{1}{2}
\begin{pmatrix}
\rho_0 & f_2 \sqrt{\rho_0}\sqrt{\rho_1}\\
f_2 \sqrt{\rho_1}\sqrt{\rho_0} & \rho_1
\end{pmatrix}
= (\cZ_p \otimes \id)(\Psi_1)
\]
because $\cZ_p$ leaves the diagonal blocks unchanged and multiplies the off-diagonal blocks by $1 - 2p = f_2$. Setting $\ket{\phi_0} = \ket{\phi_1}$ ($f_2 = 1$, $p = 0$) exhibits $\Psi_1 = A_0 A_0^{\dagger} \succeq 0$ with $\Tr \Psi_1 = \frac{1}{2}(\Tr \rho_0 + \Tr \rho_1) = 1$, so $\Psi_1 \in \cD(\bC^2 \otimes \cH_1)$. Moreover, $A_0^{\dagger} A_0 = \frac{\rho_0 + \rho_1}{2} \otimes \ketbra{\phi_0}$ shares the same nonzero eigenvalues (with multiplicities) as $A_0 A_0^{\dagger} = \Psi_1$, whence $S(\Psi_1) = S(\frac{\rho_0 + \rho_1}{2})$.

For general $f_2 \in [0,1]$, $A^{\dagger}A$ and $AA^{\dagger} = (\cZ_p \otimes \id)(\Psi_1)$ likewise have the same nonzero eigenvalues (with multiplicities) and therefore the same von Neumann entropy. Finally, $(\cZ_{1/2} \otimes \id)(\Psi_1) = \frac{1}{2}\ketbra{0} \otimes \rho_0 + \frac{1}{2}\ketbra{1} \otimes \rho_1$ is a classical--quantum block-diagonal state with entropy $1 + \frac{1}{2}(S(\rho_0) + S(\rho_1))$, so
\[
C(0) = S\bigl((\cZ_{1/2} \otimes \id)(\Psi_1)\bigr) - S(\Psi_1) = 1 - \QJS(\rho_0, \rho_1) = H(W_1). \qedhere
\]
\end{proof}

Combining Lemma~\ref{lem:gram-dephasing} with Theorem~\ref{thm:dephasing-sdpi}, Renes' channel duality (Proposition~\ref{prop:duality}), and classical posterior conditioning now proves Theorem~\ref{thm:qmgl}(a).

\begin{proof}[Proof of Theorem~\ref{thm:qmgl}(a)]
Suppose first that $W_2 : x \mapsto \ketbra{\phi_x}$ is a pure-state channel with overlap $f_2 := \lvert \braket{\phi_0}{\phi_1} \rvert = 1 - 2p$ ($p \in [0, 1/2]$), so that $H_2 = 1 - h_2(p)$ and $h_2^{-1}(1 - H_2) = p$. Adjusting the global phase of $\ket{\phi_1}$ so that $\braket{\phi_0}{\phi_1} = f_2 \ge 0$, there exists a unitary operator $U$ on $\cH_2$ exchanging $\ket{\phi_0}$ and $\ket{\phi_1}$: if $\ket{\phi_0} = \ket{\phi_1}$ we take $U = I$, and otherwise we take reflection across the bisector $\ket{\phi_0} + \ket{\phi_1}$ on $\mathrm{span}\{\ket{\phi_0}, \ket{\phi_1}\}$ extended by the identity on its orthogonal complement. Conjugation by $I \otimes U$ maps the state pair $(\rho_1 \otimes \ketbra{\phi_0}, \rho_0 \otimes \ketbra{\phi_1})$ to $(\rho_1 \otimes \ketbra{\phi_1}, \rho_0 \otimes \ketbra{\phi_0})$. By unitary invariance and symmetry of $\QJS$, the two $U_1$-branches in~\eqref{eq:var-node-entropy} coincide, and Lemma~\ref{lem:gram-dephasing} yields
\begin{align*}
H(W_1 \circledast W_2) &= 1 - \QJS\bigl(\rho_0 \otimes \ketbra{\phi_0},\; \rho_1 \otimes \ketbra{\phi_1}\bigr)\\
&= 1 + \frac{1}{2}\bigl(S(\rho_0) + S(\rho_1)\bigr) - S\!\left(\frac{\rho_0 \otimes \ketbra{\phi_0} + \rho_1 \otimes \ketbra{\phi_1}}{2}\right)\\
&= S\bigl((\cZ_{1/2} \otimes \id)(\Psi_1)\bigr) - S\bigl((\cZ_p \otimes \id)(\Psi_1)\bigr) = C(p).
\end{align*}
Applying Theorem~\ref{thm:dephasing-sdpi} with $C(0) = H_1$ gives $H(W_1 \circledast W_2) \le 1 - h_2(h_2^{-1}(1 - H_1) * p) = 1 - g_c(1 - H_1, 1 - H_2)$. By conservation of entropy~\eqref{eq:entropy-conservation},
\[
H(W_1 \boxtimes W_2) = H_1 + H_2 - H(W_1 \circledast W_2) \;\ge\; H_1 + H_2 - 1 + g_c(1 - H_1, 1 - H_2) \;\ge\; g_q(H_1, H_2),
\]
proving~\eqref{eq:qmgl-pure}.

Now suppose that $W_2$ is classical, meaning that its output states $\sigma_0, \sigma_1 \in \cD(\cH_2)$ commute and therefore correspond to conditional probability distributions on a finite alphabet $\cY$ of cardinality $\dim \cH_2$. If $W_2 = \BSC(q)$ for some $q \in [0, 1/2]$, then $H_2 = h_2(q)$ and, by Proposition~\ref{prop:duality}, $W_2^{\perp} = \Pure(1 - 2q)$ with $H(W_2^{\perp}) = 1 - H_2$ and $H(W_1^{\perp}) = 1 - H_1$. Applying~\eqref{eq:qmgl-pure} to the dual pair $(W_1^{\perp}, W_2^{\perp})$ and substituting into~\eqref{eq:duality-check} gives
\begin{align*}
H\bigl(W_1 \boxtimes \BSC(q)\bigr) &= H_1 + H_2 - 1 + H(W_1^{\perp} \boxtimes W_2^{\perp})\\
&\ge H_1 + H_2 - 1 + (1 - H_1) + (1 - H_2) - 1 + g_c(H_1, H_2) = g_c(H_1, H_2).
\end{align*}

For a general (not necessarily symmetric) classical binary-input channel $W_2 : \{0,1\} \to \cP(\cY)$ with output random variable $Y$ under uniform input $X_2$, let $r_y := \frac{1}{2}(W_2(y \mid 0) + W_2(y \mid 1))$ for $y \in \cY$. Whenever $r_y > 0$, define the posterior probability $q_y := \bP(X_2 = 1 \mid Y = y) = \frac{W_2(y \mid 1)}{W_2(y \mid 0) + W_2(y \mid 1)}$ and set $\tilde{q}_y := \min(q_y, 1 - q_y) \in [0, 1/2]$, so that $\sum_y r_y h_2(\tilde{q}_y) = H(X_2 \mid Y) = H_2$. Conditioned on $Y = y$, the bit $X_2$ is independent of $(X_1, B_1)$ with $\bP(X_2 = 1 \mid Y = y) = q_y$; since $X_1$ is uniform, $U_1 = X_1 \oplus X_2$ remains uniform and the conditional joint state of $(U_1, B_1)$ given $Y = y$ is
\[
\frac{1}{2}\ketbra{0} \otimes \bigl((1 - q_y)\rho_0 + q_y \rho_1\bigr) + \frac{1}{2}\ketbra{1} \otimes \bigl((1 - q_y)\rho_1 + q_y \rho_0\bigr).
\]
Up to relabeling $u_1 \mapsto u_1 \oplus 1$ when $q_y > 1/2$, this is the c-q state of the pre-composed channel $W_1 \circ \BSC(\tilde{q}_y)$, which coincides (up to the same input relabeling) with both classical output branches of $W_1 \boxtimes \BSC(\tilde{q}_y)$. Hence $H(X_1 \oplus X_2 \mid B_1, Y = y) = H(W_1 \circ \BSC(\tilde{q}_y)) = H(W_1 \boxtimes \BSC(\tilde{q}_y))$. By classical conditioning on $Y$ and the convexity of $v \mapsto g_c(H_1, v)$ on $[0,1]$~\cite{WynerZiv1973}, Jensen's inequality gives
\[
H(W_1 \boxtimes W_2) = \sum_y r_y H\bigl(W_1 \boxtimes \BSC(\tilde{q}_y)\bigr) \;\ge\; \sum_y r_y g_c\bigl(H_1, h_2(\tilde{q}_y)\bigr) \;\ge\; g_c(H_1, H_2) \;\ge\; g_q(H_1, H_2),
\]
proving~\eqref{eq:qmgl-classical}.
\end{proof}

\begin{remark}
When $W_1$ and $W_2$ are both classical, \eqref{eq:qmgl-classical} recovers Wyner and Ziv's Mrs.\ Gerber's Lemma~\cite{WynerZiv1973}, with equality on binary symmetric channels. When $W_1$ and $W_2$ are both pure-state channels, \eqref{eq:qmgl-pure} holds with equality.
\end{remark}

\subsection{Linear-Scaling Bound for Arbitrary Classical--Quantum Channels}

We now prove Theorem~\ref{thm:qmgl}(b) in three steps: the four-branch purification bound~\eqref{eq:four-branch} (Proposition~\ref{prop:four-branch}), the linear splitting bound~\eqref{eq:linear-splitting} together with its endpoint asymptotic~\eqref{eq:endpoint-scaling} (Proposition~\ref{prop:linear-splitting}), and the uniform potential contraction~\eqref{eq:potential-contraction} (Proposition~\ref{prop:potential-contraction}).

\begin{proposition}[Four-Branch Purification Bound]\label{prop:four-branch}
Inequality~\eqref{eq:four-branch} holds for all independent binary-input c-q channels $W_1, W_2$.
\end{proposition}

\begin{proof}
Let $W_1 : x \mapsto \rho_x \in \cD(\cH_1)$ and $W_2 : x \mapsto \sigma_x \in \cD(\cH_2)$, and set $f_2 := F(\sigma_0, \sigma_1)$. By Uhlmann's theorem~\cite{Uhlmann1976}, there exist purifications $\ket{\phi_0}, \ket{\phi_1} \in \cH_2 \otimes \cH_E$ of $\sigma_0, \sigma_1$ such that $\lvert \braket{\phi_0}{\phi_1} \rvert = f_2$. Letting $W_2^{\mathrm{pure}} : x \mapsto \ketbra{\phi_x}$, we have $W_2 = \Tr_E \circ W_2^{\mathrm{pure}}$ and $H(W_2^{\mathrm{pure}}) = 1 - h_2(\frac{1 - f_2}{2})$. Since $(W_1 \boxtimes W_2)(u_1) = (\id_{B_1} \otimes \Tr_E)((W_1 \boxtimes W_2^{\mathrm{pure}})(u_1))$, the data processing inequality for conditional von Neumann entropy followed by Theorem~\ref{thm:qmgl}(a) yields
\begin{align*}
H(W_1 \boxtimes W_2) &\ge H(W_1 \boxtimes W_2^{\mathrm{pure}}) \ge H_1 + \left(1 - h_2\!\left(\frac{1 - f_2}{2}\right)\right) - 1 + h_2\!\left(h_2^{-1}(1 - H_1) * \frac{1 - f_2}{2}\right)\\
&= H_1 + \Delta\!\left(1 - H_1, \frac{1 - f_2}{2}\right).
\end{align*}
For fixed $u \in [0,1]$ and $p_1 := h_2^{-1}(u) \in [0, 1/2]$, differentiating $\Delta(u,p) = h_2(p_1 * p) - h_2(p)$ with respect to $p \in (0, 1/2)$ gives
\[
\frac{\partial}{\partial p}\Delta(u,p) = (1 - 2p_1) h_2'(p_1 * p) - h_2'(p) \;\le\; 0,
\]
because $p \le p_1 * p \le 1/2$, $0 \le 1 - 2p_1 \le 1$, and $h_2'(s) = \log_2(\frac{1-s}{s})$ is nonnegative and decreasing on $(0, 1/2]$. Hence $p \mapsto \Delta(1 - H_1, p)$ is nonincreasing on $[0, 1/2]$.

By the fidelity--entropy inequality $f_2 \ge 2^{H_2} - 1$~\cite[Theorem~V.3]{HircheReeb2018}, we have $\frac{1 - f_2}{2} \le 1 - 2^{H_2 - 1} \in [0, 1/2]$, and therefore
\[
H(W_1 \boxtimes W_2) \;\ge\; H_1 + \Delta\bigl(1 - H_1, 1 - 2^{H_2 - 1}\bigr).
\]
Interchanging the roles of $W_1$ and $W_2$ yields the second branch of~\eqref{eq:four-branch}. Applying these first two branches to the dual pair $(W_1^{\perp}, W_2^{\perp})$ (using $H(W_i^{\perp}) = 1 - H_i$ from Proposition~\ref{prop:duality}) and substituting into~\eqref{eq:duality-check} gives the third and fourth branches:
\[
H(W_1 \boxtimes W_2) = H_1 + H_2 - 1 + H(W_1^{\perp} \boxtimes W_2^{\perp}) \;\ge\; H_2 + \Delta\bigl(H_1, 1 - 2^{-H_2}\bigr),
\]
and symmetrically $H_1 + \Delta(H_2, 1 - 2^{-H_1})$.
\end{proof}

\begin{proposition}[Linear Splitting Constant and Endpoint Asymptotics]\label{prop:linear-splitting}
Inequalities~\eqref{eq:linear-splitting} and~\eqref{eq:endpoint-scaling} hold for every binary-input c-q channel $W$.
\end{proposition}

\begin{proof}
Specialize Proposition~\ref{prop:four-branch} to $W_1 = W_2 = W$ with $H = H(W)$. The endpoints $H \in \{0,1\}$ are immediate. By the duality symmetry $H(W^-) - H(W) = H((W^{\perp})^-) - H(W^{\perp})$ from~\eqref{eq:duality-check}, it suffices to prove~\eqref{eq:linear-splitting} for $H \in [1/2, 1)$, where Proposition~\ref{prop:four-branch} gives
\[
H(W^-) - H(W) \;\ge\; \alpha(H) := \Delta(1 - H, p_2),
\]
with $\delta := 1 - H \in (0, 1/2]$, $p_1 := h_2^{-1}(\delta) \in (0, 1/2]$, $p_2 := 1 - 2^{-\delta} \in (0, 1 - 2^{-1/2}]$, and $q := p_1 * p_2 = p_2 + p_1(1 - 2p_2) \in (p_2, 1/2]$. We treat the interior and boundary regimes separately.

\emph{Region~1: $H \in [1/2, 3/4]$ ($\delta \in [1/4, 1/2]$).} Express $q$ as the convex combination $q = (1 - 2p_1)p_2 + (2p_1)\cdot \frac{1}{2}$. By concavity of $h_2$ on $[0, 1/2]$ and $h_2(1/2) = 1$,
\[
\alpha(H) = h_2(q) - h_2(p_2) \;\ge\; 2p_1\bigl(1 - h_2(p_2)\bigr).
\]
Because $h_2$ is concave and increasing on $[0, 1/2]$, its inverse $h_2^{-1}$ is convex on $[0,1]$ with $h_2^{-1}(0) = 0$; hence $\delta \mapsto h_2^{-1}(\delta)/\delta$ is nondecreasing on $(0,1]$, so for $\delta \ge 1/4$ we have $p_1 = h_2^{-1}(\delta) \ge 4h_2^{-1}(1/4)\,\delta$. Moreover, $p_2 \le 1 - 2^{-1/2}$ implies $1 - h_2(p_2) \ge 1 - h_2(1 - 2^{-1/2})$, while $H \le 3/4$ implies $\delta = 1 - H \ge \frac{4}{3} H(1 - H)$. Combining these three estimates yields
\[
\alpha(H) \;\ge\; \frac{32}{3}\, h_2^{-1}(1/4)\bigl(1 - h_2(1 - 2^{-1/2})\bigr) H(1 - H) = \theta_0\, H(1 - H),
\]
where $\theta_0 = 0.05673348\ldots > 0.05673$.

\emph{Region~2: $H \in [3/4, 1)$ ($\delta \in (0, 1/4]$).} Set $L := \log_2(1/\delta) \ge 2$. We first show that $p_1 \ge \frac{\delta}{4L}$. Let $t := \frac{\delta}{4L} \in (0, 1/32]$. Since $-(1 - t)\ln(1 - t) \le t$ for all $t \in [0, 1)$, we have $h_2(t) \le t\log_2(e/t)$, and therefore
\[
\frac{h_2(t)}{\delta} \;\le\; \frac{1}{4L}\log_2\!\left(\frac{4eL}{\delta}\right) = \frac{1}{4} + \frac{\log_2(4eL)}{4L}.
\]
The function $L \mapsto \frac{\ln(4eL)}{L}$ has derivative $\frac{1 - \ln(4eL)}{L^2} < 0$ for $L \ge 1$, so for every $L \ge 2$,
\[
\frac{h_2(t)}{\delta} \;\le\; \frac{1}{4} + \frac{\log_2(8e)}{8} = 0.8053\ldots < 1.
\]
Since $h_2$ is increasing on $[0, 1/2]$, this yields $p_1 = h_2^{-1}(\delta) \ge \frac{\delta}{4L}$.

Next, concavity of $h_2$ on $[0, 1/8]$ together with $h_2(0) = 0$ and $h_2(1/8) = 0.5435\ldots > 1/4$ implies $h_2(\delta/2) \ge \delta$ for $\delta \in (0, 1/4]$, whence $p_1 \le \delta/2$. Combined with $p_2 = 1 - e^{-(\ln 2)\delta} \le (\ln 2)\delta$, this gives
\[
q = p_1 + p_2 - 2p_1 p_2 \;\le\; \left(\frac{1}{2} + \ln 2\right)\delta < \frac{6}{5}\delta \le \frac{3}{10} < \frac{1}{2}.
\]
Consequently $\frac{1-q}{q} \ge \frac{7/10}{(6/5)\delta} = \frac{7}{12\delta}$, and for all $L \ge 2$,
\[
\log_2\!\left(\frac{1-q}{q}\right) \;\ge\; L - \log_2\!\left(\frac{12}{7}\right) \;\ge\; \left(1 - \frac{1}{2}\log_2\!\left(\frac{12}{7}\right)\right) L.
\]
By concavity of $h_2$, we have $\alpha(H) = h_2(q) - h_2(p_2) \ge (q - p_2) h_2'(q) = p_1(1 - 2p_2)\log_2(\frac{1-q}{q})$. Substituting $1 - 2p_2 = 2^H - 1 \ge 2^{3/4} - 1$, $p_1 \ge \frac{\delta}{4L}$, and $H(1 - H) \le \delta$, the logarithmic factor $L$ cancels:
\[
\alpha(H) \;\ge\; \frac{2^{3/4} - 1}{4}\left(1 - \frac{1}{2}\log_2\!\left(\frac{12}{7}\right)\right)\delta \;\ge\; \theta_2\, H(1 - H),
\]
where $\theta_2 = 0.104177\ldots > \theta_0$. Combining Regions~1 and~2 establishes~\eqref{eq:linear-splitting} for all $H \in [0,1]$.

For~\eqref{eq:endpoint-scaling}, as $\delta \to 0^+$ we have $p_1 = \frac{\delta}{\log_2(1/\delta)}(1 + o(1))$, $p_2 = (\ln 2)\delta(1 + o(1))$, and $q = (\ln 2)\delta(1 + o(1))$. Thus $q - p_2 = \frac{\delta}{\log_2(1/\delta)}(1 + o(1))$, while both $h_2'(p_2)$ and $h_2'(q)$ equal $\log_2(1/\delta)(1 + o(1))$. By concavity of $h_2$,
\[
(q - p_2) h_2'(q) \;\le\; \Delta(\delta, 1 - 2^{-\delta}) \;\le\; (q - p_2) h_2'(p_2),
\]
and both sides equal $(1 + o(1))\delta$, yielding~\eqref{eq:endpoint-scaling}.
\end{proof}

Next, we convert the linear entropy-splitting bound~\eqref{eq:linear-splitting} into a uniform multiplicative contraction of the polarization potential $T(W) = \sqrt{H(W)(1 - H(W))}$, mirroring the potential-function method used in classical finite-blocklength polarization~\cite{HassaniAlishahiUrbanke2014,GuruswamiXia2015}.

\begin{proposition}[Uniform Contraction of the Polarization Potential]\label{prop:potential-contraction}
If a binary-input c-q channel $W$ satisfies~\eqref{eq:linear-splitting}, then $T(W) := \sqrt{H(W)(1 - H(W))}$ satisfies~\eqref{eq:potential-contraction}.
\end{proposition}

\begin{proof}
If $H(W) \in \{0,1\}$, both sides of~\eqref{eq:potential-contraction} vanish. Assume $H := H(W) \in (0,1)$ and let $\alpha := H(W^-) - H = H - H(W^+) \ge \theta_0 H(1 - H)$. Because $x \mapsto \sqrt{x(1 - x)}$ is strictly concave on $[0,1]$, the symmetric sum $u \mapsto \sqrt{(H - u)(1 - H + u)} + \sqrt{(H + u)(1 - H - u)}$ is strictly decreasing for $u \in [0, \min(H, 1 - H)]$. Evaluating at $u = \theta_0 H(1 - H)$ and factoring out $T(W) = \sqrt{H(1 - H)}$ gives
\[
\frac{T(W^+) + T(W^-)}{2 T(W)} \;\le\; \frac{\sqrt{1 - a - v} + \sqrt{1 + a - v}}{2} = \sqrt{\frac{1 - v + \sqrt{(1 - v)^2 - a^2}}{2}},
\]
where $a := \theta_0(1 - 2H)$ and $v := \theta_0^2 H(1 - H)$. Since $0 < \theta_0 < 1$, setting $s := \lvert 1 - 2H \rvert \in [0,1]$ yields $1 - v - \lvert a \rvert = 1 - \frac{\theta_0^2}{4}(1 - s^2) - \theta_0 s \ge 1 - \theta_0 > 0$, so $\lvert a \rvert < 1 - v \le 1$. By the elementary inequality $\sqrt{1 - r} \le 1 - r/2$ for $r \in [0,1]$, we have $\sqrt{(1 - v)^2 - a^2} \le (1 - v) - \frac{a^2}{2(1 - v)} \le 1 - v - \frac{a^2}{2}$. Using the identity
\[
v + \frac{a^2}{4} = \theta_0^2 H(1 - H) + \frac{\theta_0^2\bigl(1 - 4H(1 - H)\bigr)}{4} = \frac{\theta_0^2}{4},
\]
we conclude that $\frac{T(W^+) + T(W^-)}{2 T(W)} \le \sqrt{1 - \frac{\theta_0^2}{4}} = \Lambda_0 < 0.99960$.
\end{proof}

\begin{proof}[Proof of Theorem~\ref{thm:qmgl}(b)]
Combine Propositions~\ref{prop:four-branch}, \ref{prop:linear-splitting}, and~\ref{prop:potential-contraction}.
\end{proof}

\begin{remark}
Numerical evaluation of $\Delta(1 - H, 1 - 2^{H-1})/(H(1 - H))$ on $[1/2, 1)$ indicates that its minimum is attained at $H = 1/2$, with $\theta_q := 4\Delta(1/2, 1 - 2^{-1/2}) = 0.20366416\ldots$ and $\Lambda_q := \sqrt{1 - \theta_q^2/4} = 0.99480160\ldots$. Table~\ref{tab:comparison} compares $\Delta(1 - H, 1 - 2^{H-1})$ against Hirche and Reeb's conjectured optimal bound $g_q(H,H) - H$~\cite[Conjecture~VII.1]{HircheReeb2018}.
\end{remark}

\begin{table}[!htbp]
\centering
\begin{tabular}{ccccc}
\toprule
$H$ & $1 - H$ & $\Delta(1 - H, 1 - 2^{H-1})$ & $g_q(H,H) - H$ & $\Delta/(1 - H)$ \\
\midrule
$0.500000$ & $0.5$ & $0.05091604$ & $0.21353673$ & $0.101832$ \\
$0.700000$ & $0.3$ & $0.06524558$ & $0.17155575$ & $0.217485$ \\
$0.900000$ & $0.1$ & $0.04135457$ & $0.07201271$ & $0.413546$ \\
$0.990000$ & $10^{-2}$ & $0.00600769$ & $0.00826493$ & $0.600769$ \\
$0.999900$ & $10^{-4}$ & $0.00007252$ & $0.00008946$ & $0.725159$ \\
$0.999999$ & $10^{-6}$ & $0.00000078$ & $0.00000092$ & $0.783046$ \\
\bottomrule
\end{tabular}
\caption{Comparison of the lower bound $\Delta(1 - H, 1 - 2^{H-1})$ from Theorem~\ref{thm:qmgl} with Hirche and Reeb's conjectured optimal bound $g_q(H,H) - H$ across $H \in [1/2, 1)$.}
\label{tab:comparison}
\end{table}

\section{Polynomial-Blocklength Classical--Quantum Polar Codes}\label{sec:polar-codes}

We now apply the uniform potential contraction~\eqref{eq:potential-contraction} to prove Theorem~\ref{thm:main-coding}. Following the standard stochastic-process formulation of channel polarization~\cite{Arikan2009,WildeGuha2013,GuruswamiXia2015}, fix a binary-input c-q channel $W_0 = W : \{0,1\} \to \cD(\cH)$ and let $(B_t)_{t \ge 1}$ be i.i.d.\ $\mathrm{Bernoulli}(1/2)$ random variables with natural filtration $\cF_t := \sigma(B_1, \ldots, B_t)$. Define the synthetic channel process inductively by $W_t := W_{t-1}^+$ if $B_t = 1$ and $W_t := W_{t-1}^-$ if $B_t = 0$. After $n$ steps, the law of $W_n$ is uniform over the $N = 2^n$ synthesized channels $\{W_n^{(i)}\}_{i=1}^{N}$ generated by Ar{\i}kan's polar transform $G_N$.

\subsection{Quantum Bhattacharyya Evolution}

To pass from conditional entropy polarization to an upper bound on the quantum successive-cancellation decoding error probability, we record the relationship between $Z(W) = F(\rho_0, \rho_1)$ and $H(W)$ together with the one-step evolution of $Z$ under $W \mapsto (W^+, W^-)$ (cf.~\cite[Proposition~5]{WildeGuha2013}).

\begin{lemma}[Quantum Bhattacharyya Evolution and Entropy Relation]\label{lem:bhattacharyya}
For every binary-input c-q channel $W : x \mapsto \rho_x \in \cD(\cH)$, the fidelity parameter $Z(W) = F(\rho_0, \rho_1)$ satisfies:
\begin{enumerate}
\item[\textup{(a)}] $Z(W) \le \sqrt{(2 \ln 2)\, H(W)}$;
\item[\textup{(b)}] $Z(W^+) = Z(W)^2$ and $Z(W^-) \le 2 Z(W)$.
\end{enumerate}
\end{lemma}

\begin{proof}
(a) By the Roga--Fannes--\.{Z}yczkowski inequality~\cite{RogaFannesZyczkowski2010} (see also~\cite[Theorem~V.3]{HircheReeb2018}), the Holevo information of a uniform binary ensemble satisfies $1 - H(W) = \QJS(\rho_0, \rho_1) \le h_2(\frac{1 - Z(W)}{2})$. Expanding $1 - h_2(\frac{1 - z}{2})$ in a power series around $z = 0$ gives
\[
1 - h_2\!\left(\frac{1 - z}{2}\right) = \frac{1}{\ln 2}\sum_{k=1}^{\infty} \frac{z^{2k}}{2k(2k - 1)} \;\ge\; \frac{z^2}{2\ln 2}
\]
for all $z \in [0,1]$, whence $H(W) \ge \frac{Z(W)^2}{2\ln 2}$.

(b) By Definition~\ref{def:polar-combinations}, the two output states of $W^+$ are block-diagonal across the classical register $U_1$:
\[
W^+(0) = \frac{1}{2}\ketbra{0} \otimes \rho_0 \otimes \rho_0 + \frac{1}{2}\ketbra{1} \otimes \rho_1 \otimes \rho_0, \qquad
W^+(1) = \frac{1}{2}\ketbra{0} \otimes \rho_1 \otimes \rho_1 + \frac{1}{2}\ketbra{1} \otimes \rho_0 \otimes \rho_1.
\]
By additivity of root fidelity over orthogonal classical blocks and multiplicativity under tensor products,
\[
Z(W^+) = \frac{1}{2} F(\rho_0 \otimes \rho_0, \rho_1 \otimes \rho_1) + \frac{1}{2} F(\rho_1 \otimes \rho_0, \rho_0 \otimes \rho_1) = F(\rho_0, \rho_1)^2 = Z(W)^2.
\]

For $W^-$, write the two output states as $\tau_0 = X_0 + X_1$ and $\tau_1 = Y_0 + Y_1$, where $X_0 := \frac{1}{2}\rho_0 \otimes \rho_0$, $X_1 := \frac{1}{2}\rho_1 \otimes \rho_1$, $Y_0 := \frac{1}{2}\rho_0 \otimes \rho_1$, and $Y_1 := \frac{1}{2}\rho_1 \otimes \rho_0$. Define block column operators $P, Q : \cH^{\otimes 2} \to \bC^2 \otimes \cH^{\otimes 2}$ by
\[
P := \begin{pmatrix} \sqrt{X_0} \\ \sqrt{X_1} \end{pmatrix}, \qquad
Q := \begin{pmatrix} \sqrt{Y_0} \\ \sqrt{Y_1} \end{pmatrix}.
\]
Since $P^{\dagger} P = \tau_0$ and $Q^{\dagger} Q = \tau_1$, polar decomposition yields isometries $U, V : \cH^{\otimes 2} \to \bC^2 \otimes \cH^{\otimes 2}$ such that $P = U\sqrt{\tau_0}$ and $Q = V\sqrt{\tau_1}$. By isometric invariance and the triangle inequality for the trace norm,
\[
Z(W^-) = \lVert \sqrt{\tau_0}\sqrt{\tau_1} \rVert_1 = \lVert P Q^{\dagger} \rVert_1 =
\left\lVert
\begin{pmatrix}
\sqrt{X_0}\sqrt{Y_0} & \sqrt{X_0}\sqrt{Y_1}\\
\sqrt{X_1}\sqrt{Y_0} & \sqrt{X_1}\sqrt{Y_1}
\end{pmatrix}
\right\rVert_1
\;\le\; \sum_{i,j \in \{0,1\}} \lVert \sqrt{X_i}\sqrt{Y_j} \rVert_1.
\]
By multiplicativity of root fidelity under tensor products, each of the four summands equals $\frac{1}{2} F(\rho_0, \rho_1) = \frac{1}{2} Z(W)$, giving $Z(W^-) \le 2 Z(W)$.
\end{proof}

\subsection{Three-Stage Polarization and Proof of Theorem~\ref{thm:main-coding}}

Fix $\varepsilon \in (0,1)$. Following the multi-stage strategy of classical finite-blocklength polarization~\cite{GuruswamiXia2015,HassaniAlishahiUrbanke2014,MondelliHassaniUrbanke2016}, we obtain a polynomial blocklength $N = 2^n = O((1/\varepsilon)^{\mu})$ without invoking the martingale convergence theorem by evolving the synthetic channel process $(W_t)_{t \ge 0}$ through three consecutive stages of depths $m_1, m_2, m_3$ (with total depth $n = m_1 + m_2 + m_3$):
\begin{itemize}
\item \emph{Stage~1 (rough polarization):} Using the uniform potential contraction~\eqref{eq:potential-contraction}, $m_1 = O(\log(1/\varepsilon))$ steps drive $H(W_{m_1})$ below $\eta = O(\varepsilon^4)$ on an event of probability at least $I(W) - \varepsilon/4$, which via Lemma~\ref{lem:bhattacharyya}(a) places $Z(W_{m_1})$ inside the quadratic-contraction regime $Z < 1/16$.
\item \emph{Stage~2 (bootstrap polarization):} Conditioning on Stage~1, a stopped $\sqrt{Z}$-supermartingale drives $Z(W_{m_1 + m_2})$ below $2^{-Y_0}$ (with $Y_0 = \Theta(\log(1/\varepsilon))$) in $m_2 = O(\log(1/\varepsilon))$ steps while losing at most $\varepsilon/4$ in probability.
\item \emph{Stage~3 (fine polarization):} Conditioning on Stage~2, a golden-ratio exponential supermartingale ensures that $\log_2(1/Z)$ never drops below $200$, while Hoeffding's inequality~\cite{Hoeffding1963} guarantees that at least a $0.499$ fraction of the $m_3$ steps are $+$ steps, driving $Z(W_n)$ below $2^{-200 N^{0.49047}}$.
\end{itemize}

\begin{lemma}[Stage 1: Rough Polarization]\label{lem:stage1}
For $\varepsilon \in (0,1)$, set
\[
m_1 := \left\lceil \frac{8}{\log_2(1/\Lambda_0)} \log_2\!\left(\frac{64}{\varepsilon}\right) \right\rceil, \qquad
\eta := \Lambda_0^{m_1/2}, \qquad
\cE_1 := \bigl\{H(W_{m_1}) \le \eta\bigr\}.
\]
Then $\bP(\cE_1) \ge I(W) - \frac{\varepsilon}{4}$. Moreover, on $\cE_1$ one has $Z(W_{m_1}) \le Z_0 := \sqrt{(2\ln 2)\eta} < \frac{1}{16}$ and $4\sqrt{Z_0} < \frac{\varepsilon}{8}$.
\end{lemma}

\begin{proof}
Iterating the conditional contraction $\bE[T(W_t) \mid \cF_{t-1}] \le \Lambda_0 T(W_{t-1})$ from~\eqref{eq:potential-contraction} and using $T(W_0) = \sqrt{H(W_0)(1 - H(W_0))} \le \frac{1}{2}$ yields $\bE[T(W_{m_1})] \le \Lambda_0^{m_1} T(W_0) \le \frac{1}{2} \Lambda_0^{m_1}$. By the definition of $m_1$, we have $\eta = \Lambda_0^{m_1/2} \le (\varepsilon/64)^4 < 1/4$. Whenever $H(W_{m_1}) \in (\eta, 1 - \eta)$, we have $T(W_{m_1}) \ge \sqrt{\eta(1 - \eta)} \ge \sqrt{\frac{3}{4}\eta}$, so Markov's inequality gives
\[
\bP\bigl[H(W_{m_1}) \in (\eta, 1 - \eta)\bigr] \;\le\; \frac{\bE[T(W_{m_1})]}{\sqrt{\eta(1 - \eta)}} \;\le\; \frac{\Lambda_0^{m_1}}{2\sqrt{\frac{3}{4}\eta}} = \frac{\Lambda_0^{3m_1/4}}{\sqrt{3}} < \Lambda_0^{m_1/2} = \eta.
\]
Using $\bE[1 - H(W_{m_1})] = I(W)$ (from~\eqref{eq:entropy-conservation}) and splitting the expectation according to whether $H(W_{m_1})$ lies in $[0, \eta]$, $(\eta, 1 - \eta)$, or $[1 - \eta, 1]$, we obtain
\[
I(W) = \bE[1 - H(W_{m_1})] \;\le\; \bP(\cE_1) + \bP\bigl[H(W_{m_1}) \in (\eta, 1 - \eta)\bigr] + \eta \;\le\; \bP(\cE_1) + 2\eta,
\]
and hence $\bP(\cE_1) \ge I(W) - 2\eta \ge I(W) - \frac{\varepsilon}{4}$. On $\cE_1$, Lemma~\ref{lem:bhattacharyya}(a) yields
\[
Z(W_{m_1}) \;\le\; \sqrt{(2\ln 2)\eta} = Z_0 \;\le\; \sqrt{2\ln 2}\left(\frac{\varepsilon}{64}\right)^2 < \frac{1}{16},
\]
and $4\sqrt{Z_0} \le 4(2\ln 2)^{1/4} \frac{\varepsilon}{64} = (0.0678\ldots)\varepsilon < \frac{\varepsilon}{8}$.
\end{proof}

\begin{lemma}[Stage 2: Bootstrap Polarization]\label{lem:stage2}
Let $\varphi := \frac{1 + \sqrt{5}}{2}$ be the golden ratio and $\gamma_2 := \frac{1/4 + \sqrt{2}}{2} = 0.832106\ldots < 1$. Set
\[
Y_0 := 200 + \left\lceil \frac{\ln(16/\varepsilon)}{\ln \varphi} \right\rceil, \qquad
m_2 := \left\lceil \frac{Y_0/2 + \log_2(8/\varepsilon)}{\log_2(1/\gamma_2)} \right\rceil,
\]
and $\cE_2 := \{Z(W_{m_1 + m_2}) \le 2^{-Y_0}\}$. Then for every realization $W_{m_1} = w$ in $\cE_1$,
\[
\bP(\cE_2 \mid W_{m_1} = w) \;\ge\; 1 - \frac{\varepsilon}{4}.
\]
In particular, whenever $\bP(\cE_1) > 0$, one has $\bP(\cE_2 \mid \cE_1) \ge 1 - \frac{\varepsilon}{4}$.
\end{lemma}

\begin{proof}
Fix any realization $W_{m_1} = w$ in $\cE_1$, so that $Z(w) \le Z_0 < 1/16$. Define the stopping time $\tau := \inf\{t \ge 0 : Z(W_{m_1 + t}) > 1/16\}$ and the nonnegative stopped process $M_t := \sqrt{Z(W_{m_1 + (t \wedge \tau)})}$ for $t \ge 0$. On the event $\{t \ge \tau\}$, we have $M_{t+1} = M_t$. On the event $\{t < \tau\}$, we have $Z(W_{m_1 + t}) \le 1/16$, so Lemma~\ref{lem:bhattacharyya}(b) gives
\[
\bE[M_{t+1} \mid \cF_{m_1 + t}] = \frac{1}{2}\left(Z(W_{m_1 + t}) + \sqrt{Z(W_{m_1 + t}^-)}\right) \;\le\; \frac{\sqrt{Z(W_{m_1 + t})} + \sqrt{2}}{2}\, M_t \;\le\; \gamma_2 M_t < M_t.
\]
Thus $(M_t)_{t \ge 0}$ is a nonnegative supermartingale with $M_0 = \sqrt{Z(w)} \le \sqrt{Z_0}$. By Doob's maximal inequality for nonnegative supermartingales,
\[
\bP[\tau \le m_2 \mid W_{m_1} = w] = \bP\!\left[\max_{0 \le t \le m_2} M_t > \frac{1}{4} \;\middle\vert\; W_{m_1} = w\right] \;\le\; 4 M_0 \;\le\; 4\sqrt{Z_0} < \frac{\varepsilon}{8}.
\]
Furthermore, since $\{\tau > t + 1\} \subseteq \{\tau > t\} \in \cF_{m_1 + t}$ and $M_{t+1} \ge 0$, the tower property of conditional expectation yields
\[
\bE\bigl[\mathbf{1}_{\{\tau > t + 1\}} M_{t+1}\bigr] \;\le\; \bE\bigl[\mathbf{1}_{\{\tau > t\}} M_{t+1}\bigr] = \bE\bigl[\mathbf{1}_{\{\tau > t\}} \bE[M_{t+1} \mid \cF_{m_1 + t}]\bigr] \;\le\; \gamma_2\, \bE\bigl[\mathbf{1}_{\{\tau > t\}} M_t\bigr].
\]
By induction on $t$, $\bE[\mathbf{1}_{\{\tau > m_2\}} M_{m_2} \mid W_{m_1} = w] \le \gamma_2^{m_2} \sqrt{Z_0}$. Setting $r := 2^{-Y_0}$, Markov's inequality and $\sqrt{Z_0} < 1$ give
\[
\bP\bigl[\tau > m_2 \text{ and } Z(W_{m_1 + m_2}) > r \;\big\vert\; W_{m_1} = w\bigr] \;\le\; \frac{\gamma_2^{m_2}\sqrt{Z_0}}{\sqrt{r}} < \gamma_2^{m_2} 2^{Y_0/2} \;\le\; \frac{\varepsilon}{8}.
\]
By the union bound, $\bP(\cE_2^c \mid W_{m_1} = w) < \frac{\varepsilon}{4}$ for every $w \in \cE_1$. Averaging over $w \in \cE_1$ when $\bP(\cE_1) > 0$ gives $\bP(\cE_2 \mid \cE_1) \ge 1 - \frac{\varepsilon}{4}$.
\end{proof}

\begin{lemma}[Stage 3: Fine Polarization via Golden-Ratio Supermartingale]\label{lem:stage3}
For any integer $m_3 \ge \max\{100(m_1 + m_2), \lceil 5 \times 10^5 \ln(16/\varepsilon) \rceil\}$ and total depth $n = m_1 + m_2 + m_3$ (with $N = 2^n$), define $\cE_3 := \{Z(W_n) \le 2^{-200 N^{0.49047}}\}$. Then for every realization $W_{m_1 + m_2} = w$ in $\cE_2$,
\[
\bP(\cE_3 \mid W_{m_1 + m_2} = w) \;\ge\; 1 - \frac{\varepsilon}{8}.
\]
\end{lemma}

\begin{proof}
Fix any realization $W_{m_1 + m_2} = w$ in $\cE_2$, and set $Y_t := \log_2(1/Z(W_{m_1 + m_2 + t})) \in [0,\infty]$ for $0 \le t \le m_3$, with the conventions $\log_2(1/0) := +\infty$ and $e^{-\infty} := 0$. On $\cE_2$, we have $Y_0 \ge 200 + \frac{\ln(16/\varepsilon)}{\ln \varphi}$. By Lemma~\ref{lem:bhattacharyya}(b), $Y_{t+1} = 2 Y_t$ whenever $B_{m_1 + m_2 + t + 1} = 1$, and $Y_{t+1} \ge Y_t - 1$ whenever $B_{m_1 + m_2 + t + 1} = 0$.

Define the stopping time $\sigma := \inf\{t \ge 0 : Y_t < 200\}$. While $t < \sigma$, we have $Y_t \ge 200 \ge 2$, so every $+$ step increases $Y_t$ by at least~$2$ and every $-$ step decreases $Y_t$ by at most~$1$; hence
\[
Y_{t+1} - Y_t \;\ge\; 3 B_{m_1 + m_2 + t + 1} - 1 \qquad (t < \sigma).
\]
Set $\lambda_0 := \ln \varphi > 0$. Since $\varphi^2 = \varphi + 1$ implies $\varphi + \varphi^{-2} = 2$, the random variable $U := e^{-\lambda_0(3B - 1)}$ for $B \sim \mathrm{Bernoulli}(1/2)$ satisfies $\bE[U] = \frac{1}{2}(e^{\lambda_0} + e^{-2\lambda_0}) = 1$. Therefore $N_t := e^{-\lambda_0(Y_{t \wedge \sigma} - 200)}$ is a nonnegative supermartingale with $N_0 \le \varphi^{-(Y_0 - 200)} \le \frac{\varepsilon}{16}$. Because $\sigma \le m_3$ implies $Y_{\sigma} < 200$ and hence $N_{\sigma} > 1$, Doob's maximal inequality yields
\[
\bP[\sigma \le m_3 \mid W_{m_1 + m_2} = w] \;\le\; \bP\!\left[\max_{0 \le t \le m_3} N_t \ge 1 \;\middle\vert\; W_{m_1 + m_2} = w\right] \;\le\; N_0 \;\le\; \frac{\varepsilon}{16}.
\]

Next, let $S := \sum_{t=1}^{m_3} B_{m_1 + m_2 + t}$ count the number of $+$ steps during Stage~3. By Hoeffding's inequality~\cite{Hoeffding1963} and the lower bound $m_3 \ge 5 \times 10^5 \ln(16/\varepsilon)$,
\[
\bP[S < 0.499\, m_3 \mid W_{m_1 + m_2} = w] \;\le\; \exp\bigl(-2(0.001)^2 m_3\bigr) \;\le\; \frac{\varepsilon}{16}.
\]
By the union bound, the event $\{\sigma > m_3\} \cap \{S \ge 0.499\, m_3\}$ has conditional probability at least $1 - \frac{\varepsilon}{8}$. On this event, $Y_t \ge 200$ for all $0 \le t \le m_3$, so every $+$ step multiplies $Y_t$ by~$2$, while every $-$ step satisfies $Y_{t+1} \ge Y_t - 1 \ge (1 - \frac{1}{200}) Y_t = 2^{-\delta_{200}} Y_t$ with $\delta_{200} := -\log_2(199/200) = 0.00723156\ldots$. Consequently,
\[
Y_{m_3} \;\ge\; 200 \cdot 2^{S - \delta_{200}(m_3 - S)} \;\ge\; 200 \cdot 2^{((1 + \delta_{200})\cdot 0.499 - \delta_{200}) m_3} > 200 \cdot 2^{0.495376\, m_3},
\]
since $(1 + \delta_{200}) \cdot 0.499 - \delta_{200} = 0.49537698\ldots > 0.495376$. Finally, the hypothesis $m_3 \ge 100(m_1 + m_2)$ implies $m_3 \ge \frac{100}{101} n$, and $0.495376 \times \frac{100}{101} = 0.49047128\ldots > 0.49047$. Thus $Y_{m_3} > 200 \cdot 2^{0.49047 n} = 200 N^{0.49047}$, which gives $Z(W_n) = 2^{-Y_{m_3}} \le 2^{-200 N^{0.49047}}$.
\end{proof}

Combining Lemmas~\ref{lem:stage1}--\ref{lem:stage3} with Wilde and Guha's quantum successive-cancellation decoding bound~\cite[Proposition~4]{WildeGuha2013} now establishes Theorem~\ref{thm:main-coding}.

\begin{proof}[Proof of Theorem~\ref{thm:main-coding}]
Fix $\varepsilon \in (0,1)$, choose $m_1$ and $m_2$ as in Lemmas~\ref{lem:stage1} and~\ref{lem:stage2}, and set
\[
m_{3,\min} := \max\bigl\{100(m_1 + m_2),\; \lceil 5 \times 10^5 \ln(16/\varepsilon) \rceil\bigr\}.
\]
Each of $m_1, Y_0, m_2, m_{3,\min}$ is bounded above by an affine function of $\log_2(1/\varepsilon)$ with absolute constants depending only on $\Lambda_0, \gamma_2, \varphi$ and independent of $W$ and $\dim \cH$. Hence there exist absolute constants $\mu, C_0 \in (0,\infty)$ (with $C_0 > 2$, ensuring $N \ge 2$ automatically) such that $n_{\min} := m_1 + m_2 + m_{3,\min} \le \mu \log_2(1/\varepsilon) + \log_2 C_0$.

For any power-of-two blocklength $N = 2^n \ge C_0(1/\varepsilon)^{\mu}$, we have $n \ge n_{\min}$, so setting $m_3 := n - m_1 - m_2 \ge m_{3,\min}$ preserves all hypotheses of Lemma~\ref{lem:stage3}. Define the information set
\[
\cA := \bigl\{i \in \{1, \ldots, N\} : Z(W_n^{(i)}) \le 2^{-200 N^{0.49047}}\bigr\}.
\]
If $I(W) \le \varepsilon$, the rate bound $R \ge I(W) - \varepsilon$ holds directly with $\cA = \emptyset$ and $P_e = 0$. Otherwise $I(W) > \varepsilon$, so $\bP(\cE_1) \ge I(W) - \frac{\varepsilon}{4} > 0$, and because Lemma~\ref{lem:stage3} holds pointwise for every realization $w \in \cE_2$, we have $\bP(\cE_3 \mid \cE_1 \cap \cE_2) \ge 1 - \frac{\varepsilon}{8}$. By Lemmas~\ref{lem:stage1}--\ref{lem:stage3},
\[
\frac{\lvert \cA \rvert}{N} \;\ge\; \bP(\cE_1 \cap \cE_2 \cap \cE_3) = \bP(\cE_1)\,\bP(\cE_2 \mid \cE_1)\,\bP(\cE_3 \mid \cE_1 \cap \cE_2) \;\ge\; \left(I(W) - \frac{\varepsilon}{4}\right)\left(1 - \frac{\varepsilon}{4}\right)\left(1 - \frac{\varepsilon}{8}\right) \;\ge\; I(W) - \varepsilon.
\]

Finally, by Wilde and Guha's quantum successive-cancellation decoder error bound~\cite[Proposition~4]{WildeGuha2013} (established via Sen's noncommutative union bound~\cite{Sen2012}; note that~\cite{WildeGuha2013} uses squared fidelity $F_{\mathrm{WG}} = Z^2$, so $\sqrt{F_{\mathrm{WG}}(W_n^{(i)})} = Z(W_n^{(i)})$), averaging uniformly over the frozen vector $u_{\cA^c} \in \{0,1\}^{\lvert \cA^c \rvert}$ yields a deterministic $G_N$-coset polar code whose average block-error probability satisfies
\[
P_e(N, R) \;\le\; 2\sqrt{\sum_{i \in \cA} \frac{1}{2} Z(W_n^{(i)})} \;\le\; 2\sqrt{N \cdot 2^{-200 N^{0.49047}}} \;\le\; 2^{-N^{0.49}}
\]
for every $N \ge 2$, since $1 + \frac{1}{2}\log_2 N - 100 N^{0.49047} \le -N^{0.49}$ for all $N \ge 2$.
\end{proof}

\subsection{Consequences for Private and Quantum Polar Coding}

By Renes and Wilde's reduction~\cite{RenesWilde2014} (building on~\cite{RenesDupuisRenner2012,WildeGuha2013Degradable}), polar codes for secret-key-assisted private classical communication and entanglement-assisted quantum communication over arbitrary qubit-input quantum channels are constructed by combining binary-input c-q polar codes for the associated amplitude and phase channels $W_A$ and $W_P$, with the required assistance rate vanishing for degradable channels. Theorem~\ref{thm:main-coding} upgrades the single-block blocklength scaling of these schemes from sub-exponential to polynomial.

\begin{corollary}\label{cor:private-quantum}
The single-block polar coding schemes of Renes and Wilde~\cite{RenesWilde2014} for secret-key-assisted private classical communication and entanglement-assisted quantum communication over arbitrary finite-dimensional qubit-input quantum channels (with vanishing assistance rate for degradable channels) achieve their respective symmetric coherent or private information rates within an additive gap $\varepsilon \in (0,1)$ with decoding and secrecy error at most $O(2^{-N^{0.49}})$ at blocklength $N = O((1/\varepsilon)^{\mu})$ for a universal constant $\mu < \infty$.
\end{corollary}

\begin{proof}
Given a finite-dimensional qubit-input quantum channel $\cN$, let $W_A$ and $W_P$ denote the associated binary-input c-q amplitude and phase channels of~\cite{RenesWilde2014} (where $W_P$ is the coherent-information phase channel for quantum communication or the wiretap private-information phase channel for private classical communication), whose symmetric capacities satisfy $I(W_A) + I(W_P) - 1 = I_{\mathrm{sym}}(\cN)$. Apply the three-stage polarization analysis of Lemmas~\ref{lem:stage1}--\ref{lem:stage3} with gap parameter $\varepsilon' := \varepsilon/4$ to $W_A$ and $W_P$ simultaneously. For every power-of-two blocklength $N = 2^n \ge C_0(4/\varepsilon)^{\mu}$, the good sets $\cG_N(W_A) := \{i : Z((W_A)_n^{(i)}) \le 2^{-200 N^{0.49047}}\}$ and $\cG_N(W_P) := \{i : Z((W_P)_n^{(i)}) \le 2^{-200 N^{0.49047}}\}$ (with the phase-channel indices aligned in the dual order of~\cite{RenesWilde2014}) satisfy $\lvert \cG_N(W_A) \rvert / N \ge I(W_A) - \varepsilon/4$ and $\lvert \cG_N(W_P) \rvert / N \ge I(W_P) - \varepsilon/4$. Following~\cite{RenesWilde2014}, partition $\{1, \ldots, N\}$ into four sets: the information set $\cA := \cG_N(W_A) \cap \cG_N(W_P)$, the $Z$-frozen set $\cG_N(W_A)^c \cap \cG_N(W_P)$, the $X$-frozen set $\cG_N(W_A) \cap \cG_N(W_P)^c$, and the assistance set $\cE := \cG_N(W_A)^c \cap \cG_N(W_P)^c$ (requiring pre-shared entanglement for quantum communication or secret key for private classical communication). Then $\cA$ has rate
\[
\frac{\lvert \cA \rvert}{N} \;\ge\; \frac{\lvert \cG_N(W_A) \rvert + \lvert \cG_N(W_P) \rvert - N}{N} \;\ge\; I(W_A) + I(W_P) - 1 - \frac{\varepsilon}{2} = I_{\mathrm{sym}}(\cN) - \frac{\varepsilon}{2}.
\]
Applying the same polarization bound with $\varepsilon' = \varepsilon/4$ to the dual channels $W_A^{\perp}, W_P^{\perp}$ ensures that the complementary unpolarized fraction is at most $\varepsilon/2$; for degradable channels, degradability implies that every index in $\cE$ lies in the unpolarized set as in~\cite{RenesWilde2014}, so the required assistance rate $\lvert \cE \rvert / N \le \varepsilon/2$ vanishes as $\varepsilon \to 0$. By the successive-cancellation decoding and secrecy analysis of~\cite{RenesWilde2014} (via~\cite[Proposition~4]{WildeGuha2013}), the combined trace-distance error is at most $O(\sqrt{N \cdot 2^{-200 N^{0.49047}}}) = O(2^{-N^{0.49}})$.
\end{proof}

\begin{remark}
Unassisted quantum polar codes for arbitrary (non-degradable) qubit-input channels require two-level or chained constructions such as that of Renes, Sutter, Dupuis, and Renner~\cite{RenesSutterDupuisRenner2015}; combining~\cite{RenesSutterDupuisRenner2015} with Theorem~\ref{thm:main-coding} yields polynomial total blocklength, though amortizing pre-shared entanglement across $k = \Theta(1/\varepsilon)$ blocks increases the scaling exponent accordingly.
\end{remark}

\section{Concluding Remarks}\label{sec:concluding}

Three natural questions remain open. First, while Theorem~\ref{thm:qmgl}(a) proves Hirche and Reeb's Quantum Mrs.\ Gerber's Lemma~\cite[Conjecture~VII.1]{HircheReeb2018} whenever one of the two combined channels is pure-state or classical, establishing $H(W_1 \boxtimes W_2) \ge g_q(H_1, H_2)$ for general pairs of mixed-state noncommuting c-q channels remains open. In the proof of Proposition~\ref{prop:four-branch}, the gap relative to $g_q$ enters through replacing $W_2$ by its Uhlmann purification and bounding $F(\sigma_0, \sigma_1) \ge 2^{H_2} - 1$; eliminating this loss would likely require a bipartite generalization of the nonlinear dephasing contraction in Theorem~\ref{thm:dephasing-sdpi}.

Second, our proof of Theorem~\ref{thm:main-coding} establishes a universal finite scaling exponent $\mu < \infty$ (the parameters in Lemmas~\ref{lem:stage1}--\ref{lem:stage3} give $\mu \approx 1.39 \times 10^6$; replacing $\Lambda_0$ with $\Lambda_q \approx 0.99480$ reduces $\mu$ by a factor of roughly four to $\approx 3.48 \times 10^5$, where the Stage~3 Hoeffding floor $5 \times 10^5 \ln(16/\varepsilon)$ dominates, and relaxing the Hoeffding threshold $0.499$ slightly achieves a reduction of over an order of magnitude). Tightening the upper bound on $\mu$ toward the classical scaling-exponent benchmarks ($\mu \approx 3.627$ for the binary erasure channel, and proved general-BMS upper bounds down to $\mu \le 4.714$~\cite{HassaniAlishahiUrbanke2014,GoldinBurshtein2014,GuruswamiXia2015,MondelliHassaniUrbanke2016}), and extending the finite-blocklength analysis to non-stationary c-q channels via~\cite[Lemma~VIII.1]{HircheReeb2018} or to general non-binary alphabets~\cite{NasserRenes2018}, are natural directions for further work.

Third, Theorem~\ref{thm:main-coding} determines the blocklength scaling of classical--quantum polar codes under Wilde and Guha's quantum successive-cancellation decoder, while whether a gate-efficient quantum circuit of size $\mathrm{poly}(N)$ can implement this decoder for general c-q channels remains open.

\section*{Use of Artificial Intelligence}

The author used Gemini 4 for mathematical assistance during the preparation of this manuscript and its Lean~4 formalization, including support with intermediate derivations, verification of symbolic and numerical estimates, and formalization of the arguments.

\begin{thebibliography}{99}

\bibitem{Arikan2009}
E.~Ar{\i}kan, \emph{Channel polarization: A method for constructing capacity-achieving codes for symmetric binary-input memoryless channels}, IEEE Trans.\ Inf.\ Theory \textbf{55} (2009), no.~7, 3051--3073.

\bibitem{ArikanTelatar2009}
E.~Ar{\i}kan and E.~Telatar, \emph{On the rate of channel polarization}, in Proc.\ IEEE Int.\ Symp.\ Inf.\ Theory (ISIT), 2009, pp.~1493--1495.

\bibitem{BaumgratzCramerPlenio2014}
T.~Baumgratz, M.~Cramer, and M.~B. Plenio, \emph{Quantifying coherence}, Phys.\ Rev.\ Lett.\ \textbf{113} (2014), 140401.

\bibitem{BeigiDattaRouze2020}
S.~Beigi, N.~Datta, and C.~Rouz\'{e}, \emph{Quantum reverse hypercontractivity: Its tensorization and application to strong converses}, Comm.\ Math.\ Phys.\ \textbf{376} (2020), no.~2, 753--794.

\bibitem{BertaLamiTomamichel2024}
M.~Berta, L.~Lami, and M.~Tomamichel, \emph{Continuity of entropies via integral representations}, IEEE Trans.\ Inf.\ Theory \textbf{71} (2025), no.~3, 1896--1908.

\bibitem{FawziRenner2015}
O.~Fawzi and R.~Renner, \emph{Quantum conditional mutual information and approximate Markov chains}, Comm.\ Math.\ Phys.\ \textbf{340} (2015), no.~2, 575--611.

\bibitem{Frenkel2023}
P.~E. Frenkel, \emph{Integral formula for quantum relative entropy implies data processing inequality}, Quantum \textbf{7} (2023), 1102.

\bibitem{FuchsVanDeGraaf1999}
C.~A. Fuchs and J.~van~de~Graaf, \emph{Cryptographic distinguishability measures for quantum-mechanical states}, IEEE Trans.\ Inf.\ Theory \textbf{45} (1999), no.~4, 1216--1227.

\bibitem{Gao2015}
J.~Gao, \emph{Quantum union bounds for sequential projective measurements}, Phys.\ Rev.\ A \textbf{92} (2015), 052331.

\bibitem{GoldinBurshtein2014}
D.~Goldin and D.~Burshtein, \emph{Improved bounds on the finite length scaling of polar codes}, IEEE Trans.\ Inf.\ Theory \textbf{60} (2014), no.~11, 6966--6978.

\bibitem{GuruswamiXia2015}
V.~Guruswami and P.~Xia, \emph{Polar codes: Speed of polarization and polynomial gap to capacity}, IEEE Trans.\ Inf.\ Theory \textbf{61} (2015), no.~1, 3--16.

\bibitem{HassaniAlishahiUrbanke2014}
S.~H. Hassani, K.~Alishahi, and R.~L. Urbanke, \emph{Finite-length scaling for polar codes}, IEEE Trans.\ Inf.\ Theory \textbf{60} (2014), no.~10, 5875--5898.

\bibitem{He2026}
Z.~He, \emph{Erasure extremality and rigidity for binary information combining with quantum side information}, Zenodo record 24377765 (November 2026), \url{https://zenodo.org/records/24377765}.

\bibitem{Helstrom1976}
C.~W. Helstrom, \emph{Quantum Detection and Estimation Theory}, Academic Press, New York, 1976.

\bibitem{HircheReeb2018}
C.~Hirche and D.~Reeb, \emph{Bounds on information combining with quantum side information}, IEEE Trans.\ Inf.\ Theory \textbf{64} (2018), no.~7, 4739--4757.

\bibitem{HircheRouzeFranca2022}
C.~Hirche, C.~Rouz\'{e}, and D.~Stilck~Fran\c{c}a, \emph{On contraction coefficients, partial orders and approximation of capacities for quantum channels}, Quantum \textbf{6} (2022), 862.

\bibitem{HircheTomamichel2024}
C.~Hirche and M.~Tomamichel, \emph{Quantum R\'{e}nyi and $f$-divergences from integral representations}, Comm.\ Math.\ Phys.\ \textbf{405} (2024), no.~9, 208.

\bibitem{Hoeffding1963}
W.~Hoeffding, \emph{Probability inequalities for sums of bounded random variables}, J.\ Amer.\ Statist.\ Assoc.\ \textbf{58} (1963), no.~301, 13--30.

\bibitem{Holevo1973}
A.~S. Holevo, \emph{Bounds for the quantity of information transmitted by a quantum communication channel}, Probl.\ Peredachi Inf.\ \textbf{9} (1973), no.~3, 3--11.

\bibitem{Holevo1998}
A.~S. Holevo, \emph{The capacity of the quantum channel with general signal states}, IEEE Trans.\ Inf.\ Theory \textbf{44} (1998), no.~1, 269--273.

\bibitem{KastoryanoTemme2013}
M.~J. Kastoryano and K.~Temme, \emph{Quantum logarithmic Sobolev inequalities and rapid mixing}, J.\ Math.\ Phys.\ \textbf{54} (2013), no.~5, 052202.

\bibitem{KimRuskai2014}
I.~H. Kim and M.~B. Ruskai, \emph{Bounds on the concavity of quantum entropy}, J.\ Math.\ Phys.\ \textbf{55} (2014), no.~9, 092201.

\bibitem{LiebRuskai1973}
E.~H. Lieb and M.~B. Ruskai, \emph{Proof of the strong subadditivity of quantum-mechanical entropy}, J.\ Math.\ Phys.\ \textbf{14} (1973), no.~12, 1938--1941.

\bibitem{Lindblad1975}
G.~Lindblad, \emph{Completely positive maps and entropy inequalities}, Comm.\ Math.\ Phys.\ \textbf{40} (1975), no.~2, 147--151.

\bibitem{MondelliHassaniUrbanke2016}
M.~Mondelli, S.~H. Hassani, and R.~L. Urbanke, \emph{Unified scaling of polar codes: Error exponent, scaling exponent, moderate deviations, and error floors}, IEEE Trans.\ Inf.\ Theory \textbf{62} (2016), no.~12, 6698--6712.

\bibitem{MullerHermesReeb2017}
A.~M\"{u}ller-Hermes and D.~Reeb, \emph{Monotonicity of the quantum relative entropy under positive maps}, Ann.\ Henri Poincar\'{e} \textbf{18} (2017), no.~5, 1777--1788.

\bibitem{MullerHermesFrancaWolf2016}
A.~M\"{u}ller-Hermes, D.~Stilck~Fran\c{c}a, and M.~M. Wolf, \emph{Relative entropy convergence for depolarizing channels}, J.\ Math.\ Phys.\ \textbf{57} (2016), no.~2, 022202.

\bibitem{NasserRenes2018}
R.~Nasser and J.~M. Renes, \emph{Polar codes for arbitrary classical--quantum channels and arbitrary cq-MACs}, IEEE Trans.\ Inf.\ Theory \textbf{64} (2018), no.~11, 7424--7442.

\bibitem{OskoueiManciniWilde2019}
S.~K. Oskouei, S.~Mancini, and M.~M. Wilde, \emph{Union bound for quantum information processing}, Proc.\ R.\ Soc.\ A \textbf{475} (2019), no.~2221, 20180612.

\bibitem{PolyanskiyWu2017}
Y.~Polyanskiy and Y.~Wu, \emph{Strong data-processing inequalities for channels and Bayesian networks}, in \emph{Convexity and Concentration}, IMA Vol.\ Math.\ Appl., vol.~161, Springer, New York, 2017, pp.~211--249.

\bibitem{Raginsky2016}
M.~Raginsky, \emph{Strong data processing inequalities and $\Phi$-Sobolev inequalities for discrete channels}, IEEE Trans.\ Inf.\ Theory \textbf{62} (2016), no.~6, 3355--3389.

\bibitem{Renes2018}
J.~M. Renes, \emph{Duality of channels and codes}, IEEE Trans.\ Inf.\ Theory \textbf{64} (2018), no.~1, 577--592.

\bibitem{RenesDupuisRenner2012}
J.~M. Renes, F.~Dupuis, and R.~Renner, \emph{Efficient polar coding of quantum information}, Phys.\ Rev.\ Lett.\ \textbf{109} (2012), 050504.

\bibitem{RenesSutterDupuisRenner2015}
J.~M. Renes, D.~Sutter, F.~Dupuis, and R.~Renner, \emph{Efficient quantum polar codes requiring no preshared entanglement}, IEEE Trans.\ Inf.\ Theory \textbf{61} (2015), no.~11, 6395--6414.

\bibitem{RenesWilde2014}
J.~M. Renes and M.~M. Wilde, \emph{Polar codes for private and quantum communication over arbitrary channels}, IEEE Trans.\ Inf.\ Theory \textbf{60} (2014), no.~6, 3090--3103.

\bibitem{RogaFannesZyczkowski2010}
W.~Roga, M.~Fannes, and K.~\.{Z}yczkowski, \emph{Universal bounds for the Holevo quantity, coherent information, and the Jensen--Shannon divergence}, Phys.\ Rev.\ Lett.\ \textbf{105} (2010), 040505.

\bibitem{SasonVerdu2016}
I.~Sason and S.~Verd\'{u}, \emph{$f$-divergence inequalities}, IEEE Trans.\ Inf.\ Theory \textbf{62} (2016), no.~11, 5973--6006.

\bibitem{SchumacherWestmoreland1997}
B.~Schumacher and M.~D. Westmoreland, \emph{Sending classical information via noisy quantum channels}, Phys.\ Rev.\ A \textbf{56} (1997), no.~1, 131--138.

\bibitem{Sen2012}
P.~Sen, \emph{Achieving the Han--Kobayashi inner bound for the quantum interference channel}, in Proc.\ IEEE Int.\ Symp.\ Inf.\ Theory (ISIT), 2012, pp.~736--740.

\bibitem{Tomamichel2016}
M.~Tomamichel, \emph{Quantum Information Processing with Finite Resources: Mathematical Foundations}, SpringerBriefs in Math.\ Phys., vol.~5, Springer, Cham, 2016.

\bibitem{Uhlmann1976}
A.~Uhlmann, \emph{The ``transition probability'' in the state space of a $*$-algebra}, Rep.\ Math.\ Phys.\ \textbf{9} (1976), no.~2, 273--279.

\bibitem{Umegaki1962}
H.~Umegaki, \emph{Conditional expectation in an operator algebra, IV (entropy and information)}, Kodai Math.\ Sem.\ Rep.\ \textbf{14} (1962), no.~2, 59--85.

\bibitem{Wilde2013}
M.~M. Wilde, \emph{Quantum Information Theory}, Cambridge University Press, Cambridge, 2013.

\bibitem{WildeGuha2013}
M.~M. Wilde and S.~Guha, \emph{Polar codes for classical--quantum channels}, IEEE Trans.\ Inf.\ Theory \textbf{59} (2013), no.~2, 1175--1187.

\bibitem{WildeGuha2013Degradable}
M.~M. Wilde and S.~Guha, \emph{Polar codes for degradable quantum channels}, IEEE Trans.\ Inf.\ Theory \textbf{59} (2013), no.~7, 4718--4729.

\bibitem{WynerZiv1973}
A.~D. Wyner and J.~Ziv, \emph{A theorem on the entropy of certain binary sequences and applications---I}, IEEE Trans.\ Inf.\ Theory \textbf{19} (1973), no.~6, 769--772.

\end{thebibliography}

\bigskip
\noindent\textsc{Prince Anand}\\
\emph{E-mail address:} \texttt{princeanand123pc@gmail.com}

\end{document}
